% ===================================================================
% FST-I: Thermodynamic Stability of Fundamental Parameters
% Target: Foundations of Physics
% Version: 3.0 (source-audit revision)
% ===================================================================

\documentclass[12pt,a4paper]{article}

% Packages
\usepackage[utf8]{inputenc}
\usepackage[english]{babel}
\usepackage[T1]{fontenc}
\usepackage{amsmath,amssymb,amsthm}
\usepackage{physics}
\usepackage{graphicx}
\usepackage{booktabs}
\usepackage[margin=2.5cm]{geometry}
\usepackage{natbib}
\usepackage{enumitem}
\usepackage[expansion=false]{microtype}
\usepackage{tabularx}
\usepackage{array}
\usepackage[unicode=true]{hyperref}
\usepackage{cleveref}
\renewcommand*{\HyperDestNameFilter}[1]{en.#1}

% Theorem environments
\newtheorem{proposition}{Proposition}
\newtheorem{conjecture}{Conjecture}
\newtheorem{definition}{Definition}
\newtheorem{assumption}{Assumption}
\newtheorem{lemma}{Lemma}
\theoremstyle{remark}
\newtheorem{remark}{Remark}
\newcolumntype{Y}{>{\raggedright\arraybackslash}X}
\setlength{\emergencystretch}{2em}

% Title
\title{Functional Stability Theory I: A Game-Theoretic Framework\\for the Thermodynamic Stability of Fundamental Parameters}

\author{Lukas Geiger\\
\textit{Independent Researcher, Bernau, Germany}\\
ORCID: \url{https://orcid.org/0009-0005-7296-1534}}

\date{March 2026}

\hypersetup{
  pdftitle={Functional Stability Theory I: A Game-Theoretic Framework for the Thermodynamic Stability of Fundamental Parameters},
  pdfauthor={Lukas Geiger},
  pdfsubject={Thermodynamic stability of fundamental parameters in a game-theoretic framework},
  pdfkeywords={game theory, entropy production, Standard Model, fine-tuning, Nash equilibrium, quantum mechanics, symmetry breaking, envariance, decoherence}
}

\begin{document}

\maketitle

\begin{abstract}
This paper proposes a game-theoretic interpretive framework for the Standard Model of particle physics, grounded in thermodynamic optimization principles. We explore the hypothesis that the observed particle content and coupling constants are consistent with a Nash equilibrium of a multi-scale thermodynamic game, where the payoff function is the integrated entropy production over cosmological timescales. Within this framework, quantum mechanics is reinterpreted as a mixed-strategy equilibrium, with the Born rule recast via environment-assisted invariance (envariance) as the unique equilibrium distribution---a conceptual reinterpretation, not a new derivation. Symmetry breaking corresponds to phase transitions between equilibria, and fine-tuning is reinterpreted as constrained optimization rather than coincidence. Unlike previous thermodynamic approaches to fine-tuning that address single parameters in isolation, FST-I provides a hierarchical multi-scale framework where parameters are selected by stability across coupled sectors. The approach does not modify the mathematical formalism of quantum field theory; it provides no quantitative predictions beyond established physics but offers a unifying interpretive vocabulary grounded in stochastic game theory and mean-field dynamics. A semi-analytic stellar model demonstrates that the observed fine-structure constant lies within a broad plateau of near-maximal integrated stellar entropy production ($\mathcal{S}/\mathcal{S}_{\max} = 0.988$), providing a first quantitative consistency check---though not a decisive test---for the entropy-optimality hypothesis. The plateau itself (with $\mathcal{S}/\mathcal{S}_{\max} > 0.92$ across the entire viable range) indicates that stellar entropy production alone does not sharply select $\alpha$; a two-dimensional scan over $\alpha$ and $m_e/m_p$ yields $\mathcal{S}/\mathcal{S}_{\max}^{2D} = 0.47$, confirming that hierarchical multi-scale constraints from atomic and chemical physics are required. We establish falsifiability criteria and identify the constraints required for a complete multi-parameter test. The framework is at the stage of a structured hypothesis requiring computational validation, not a predictive theory.
\end{abstract}

\textbf{Keywords:} game theory, entropy production, Standard Model, fine-tuning, Nash equilibrium, quantum mechanics, symmetry breaking, envariance, decoherence

\clearpage

% ===================================================================
\section{Introduction: The Fine-Tuning Problem in Fundamental Physics}
\label{sec:introduction}
% ===================================================================

The pursuit of a unified theoretical framework in fundamental physics has traditionally relied on axiomatic approaches, wherein the core tenets of quantum mechanics and the parameter values of the Standard Model are accepted as irreducible givens. The historical separation between the microscopic quantum domain and macroscopic thermodynamic behavior has led to profound interpretational challenges, most notably the measurement problem and the inexplicable fine-tuning of cosmological constants \citep{Adams2019}.

The Standard Model contains 19 free parameters\footnote{Three gauge couplings ($\alpha_{EM}$, $\alpha_s$, $\alpha_W$), six quark masses, three charged lepton masses, four CKM parameters (three mixing angles and one CP-violating phase), the Higgs mass, the Higgs VEV, and the QCD vacuum angle $\theta_{\text{QCD}}$. Including neutrino masses and PMNS mixing raises the count to $\sim$26.}---coupling constants, masses, and mixing angles---whose values must be determined experimentally rather than derived from first principles. These values exhibit remarkable ``fine-tuning'': small variations would prevent nuclear binding, destabilize atoms, or fundamentally alter the mass hierarchy. This observation has prompted extensive discussion, with proposed explanations ranging from the anthropic principle \citep{Weinberg1987,Susskind2005} to landscape scenarios in string theory \citep{Bousso2000}.

The proposed framework, Functional Stability Theory~I (FST-I), offers a new interpretive perspective. It explores the hypothesis that the axiomatic structure of quantum mechanics and the highly specific fine-tuning of the Standard Model may be understood as the emergent, constrained equilibrium outcomes of a multi-scale optimization game governed by thermodynamic principles. FST-I does not derive parameter values from first principles and does not generate quantitative predictions beyond established physics; its contribution is a unifying conceptual vocabulary that organizes known parameter sensitivities within a game-theoretic architecture.

In this paradigm, the universe is modeled as a complex, non-equilibrium thermodynamic system whose large-scale dynamics are consistent with high entropy production rates over cosmological epochs \citep{MEPP2004}. The effective degrees of freedom in this game are the vacuum field sectors, whose couplings determine thermodynamic payoffs defined by local entropy production. By bridging non-cooperative game theory, environment-induced superselection (einselection), and the Maximum Entropy Production Principle (MEPP), FST-I offers a unified interpretive framework for the origin of quantum probabilities, the necessity of spontaneous symmetry breaking, and the location of physical constants in the parameter space \citep{QuantumOptimization,LasryLions2007,HiggsVEV2024}.

% ===================================================================
\section{Structural Alignment: FST-I as a Pattern A Instance}
\label{sec:structural-alignment}
% ===================================================================

The game-theoretic framing of the Standard Model provided in this work is not a mere heuristic but a realization of the \textbf{Pattern A Stability Logic} established across the FST research ecosystem (see \cite{FST-Unified2026}). Pattern~A denotes a recurring three-step structure: (i)~an infinite-dimensional configuration space is reduced to a finite effective dimension by constraint, (ii)~residual instabilities signal hidden degrees of freedom requiring gauge-fixing, and (iii)~the resulting constrained functional is positive-definite, certifying stability. The recent \textbf{gauge-theoretic reformulation of the Riemann Hypothesis} (see \cite{GeigerRH2026c}) provides a structural template for this architecture.

Within this generalized stability framework, the Standard Model equilibrium is understood as follows:
\begin{enumerate}[label=\alph*)]
    \item \textbf{Analytical Rank-Finiteness (Null-Tail):} The infinite-dimensional Hilbert space of field configurations remains tractable because the "strategically viable" subspace is analytically bounded. The high-energy (UV) tail is neutralized by the entropy-production constraint, preventing information leakage into unphysical sectors.
    \item \textbf{Finite Negativity as Gauge-Signal:} The existence of "unbroken" or massless symmetries is the analog of the finite negativity in the arithmetical kernel. These are not defects, but signals of hidden degrees of freedom that must be stabilized by a gauge-shift—represented here by symmetry breaking and the Higgs mechanism.
    \item \textbf{Constrained Functional Positivity (Nash Stability):} The Nash equilibrium (Standard Model) is the "gauge-stable" state where the free-energy functional reaches a constrained minimum. Just as the RH is reduced to a rank-2 stabilization, the Standard Model stability is the result of a finite set of parameter-couplings ensuring global thermodynamic positivity.
\end{enumerate}

This alignment transforms FST-I from a plausible analogy into a \emph{structural} analogy of the ``Functional Positivity under Constraint'' mechanism. We emphasize that this remains an analogy---the formal proof that the Standard Model parameter space admits a resolvent-stable fixed point in the sense of \citep{GeigerRH2026b} has not been established. The value of the analogy lies in organizing the known parameter sensitivities (Sections~\ref{sec:finetuning}--\ref{sec:predictions}) within a unified conceptual architecture. A first quantitative consistency check is provided by the semi-analytic stellar model in Section~\ref{sec:stellar_model}, where the observed fine-structure constant achieves 98.8\% of the maximum stellar entropy production (Table~\ref{tab:alpha_scan}).

\paragraph{Terminology.} Throughout this paper, the term \emph{resolvent-stable} denotes, by structural analogy with the spectral gap property in \citep{GeigerRH2026b}, a fixed point of the dynamics where the second-order perturbation analysis (Hessian or resolvent expansion) yields strictly positive eigenvalues in all non-gauge directions. This terminology is used as a conceptual shorthand; the formal transfer to physical systems requires domain-specific verification in each case, which has not yet been carried out for the Standard Model parameter space.

% ===================================================================
\section{Foundations of Stochastic Game Theory in Physical Systems}
\label{sec:foundations}
% ===================================================================

To comprehend the application of game theory to fundamental fields, one must transition from classical finite-agent models to continuous, stochastic mean-field games.

\subsection{From Nash Equilibrium to Mean Field Games}

In classical microeconomics and decision theory, a Nash equilibrium \citep{Nash1950} describes a state where no agent can unilaterally deviate from their chosen strategy to improve their payoff, given the strategies of all other agents. When extended to quantum systems, the classical game theory formalism is revealed to be a restricted special case of a broader quantum game theory \citep{QuantumOptimization}, where the von Neumann entropy and the von Neumann-Morgenstern utility function share a common mathematical structure (both are real-valued, concave functionals on convex state spaces), and the Hamiltonian operator can be interpreted as encoding the strategic payoff structure \citep{Abel2023,QuantumOptimization}.

As the number of interacting agents (field excitations) approaches infinity ($N \to \infty$), the discrete Nash equilibrium framework transitions into a Mean Field Game (MFG) equilibrium \citep{LasryLions2007}. In this limit, individual agents optimize their cost functionals against the empirical distribution of the entire population, driven by controlled McKean-Vlasov-type stochastic differential equations \citep{LasryLions2006}. Recent work has established refined $L^\infty$ approximation bounds for Nash equilibria in Wasserstein-space formulations of MFGs \cite{DjeteTouzi2026}, providing a rigorous metric-space framework for the convergence of finite-agent equilibria to their mean-field limits.

The measure $m(p,t)$ is interpreted as a probability distribution over parameter configurations encoding the relative stability of dynamical realizations---not as an ontological commitment to a multiverse. The MFG framework remains purely mathematical; physical content enters only through the specification of the Lagrangian~$L$ and the value function~$V$.

\subsection{Hamilton-Jacobi-Bellman Dynamics}

The evolution of these systems relies on dynamic programming and fixed-point arguments, specifically governed by the Hamilton-Jacobi-Bellman equation:
\begin{equation}
\frac{\partial V}{\partial t} + \mathcal{H}\left(x, \nabla V, t\right) = 0
\end{equation}
where the value function $V$ and the control Hamiltonian $\mathcal{H}$ dictate the optimal feedback policies \citep{LasryLions2007}. This stochastic game framework forms the mathematical bedrock for interpreting quantum states not as physical waves, but as probability distributions over a complex strategy space.

\subsection{Quantum--Strategy Correspondence}

\paragraph{Classical embedding.}
Let $\Delta_n = \{p \in \mathbb{R}_{\geq 0}^n : \sum_i p_i = 1\}$ denote the simplex of classical mixed strategies over $n$ pure strategies, and let $\mathcal{D}(\mathcal{H})$ be the set of density matrices on $\mathcal{H} = \mathbb{C}^n$. There is a canonical affine embedding $\iota: \Delta_n \hookrightarrow \mathcal{D}(\mathcal{H})$ given by $\iota(p) = \sum_i p_i |i\rangle\langle i|$, whose image is precisely the commutative (diagonal) subset of $\mathcal{D}(\mathcal{H})$. For any diagonal observable $A = \sum_i a_i |i\rangle\langle i|$, the expectation value satisfies $\mathrm{tr}(\iota(p) A) = \sum_i p_i a_i$, reproducing exactly the classical expected payoff structure.

\paragraph{Noncommutative extension.}
Density matrices thus provide a \emph{noncommutative extension} of classical mixed strategies: diagonal states reproduce classical randomization, while off-diagonal terms encode basis-dependent interference structure with no classical simplex representation. Off-diagonal elements $\rho_{ij}$ ($i \neq j$) represent phase-sensitive couplings between basis alternatives under unitary evolution and noncommuting observables---they are \emph{not} classical correlations between pure strategies \citep{Lindblad1976}.

\paragraph{Status and MFG compatibility.}
We treat this quantum--strategy correspondence as a \emph{commutative specialization}, not as an isomorphism of state spaces. Precisely: under a Lindblad optimal-transport geometry on~$\mathcal{D}(\mathcal{H})$ whose generator preserves the diagonal subalgebra and induces a Markov generator~$Q$ on the diagonal, three properties are expected to hold: (i)~diagonal initial states remain diagonal under the forward dynamics; (ii)~the noncommutative OT metric reduces on diagonals to the discrete Maas--Mielke--Erbar OT metric on~$\Delta_n$; (iii)~the coupled forward--backward system reduces to the finite-state Lasry--Lions master equation. These properties are stated as a conjecture; a rigorous proof under the stated geometric conditions has not been published and constitutes an open mathematical problem (see also Table~\ref{tab:status-claims}). If established, classical mixed strategies on~$\Delta_n$ would be a commutative special case of a quantum MFG on~$\mathcal{D}(\mathcal{H})$.

However, a genuine quantum correspondence incorporating coherences (off-diagonal elements) \emph{cannot} be recovered from classical~$W_2$-MFG alone. Any unitary component $\partial_t \rho = -i[H, \rho] + \cdots$ with non-diagonal~$H$ immediately drives the dynamics off the diagonal submanifold, and geodesics in quantum OT generically do not preserve~$\iota(\Delta_n)$. Coherences require a noncommutative extension of the state space and the analysis that goes beyond classical Lasry--Lions theory.

The rest of this paper builds on the commutative correspondence as a conceptual organizing principle and interpretive framework. The extension to coherences remains an open problem for future work.

% ===================================================================
\section{Quantum Mechanics as Mixed Strategy}
\label{sec:quantum}
% ===================================================================

To ground the thermodynamic game mathematically, we map the formal components of quantum mechanics onto the elements of game theory. In the proposed framework, the axiomatic structure of quantum mechanics is reinterpreted---at the level of vocabulary, not new physics---as a thermodynamic optimization game whose degrees of freedom are the vacuum field sectors. This mapping is a \emph{conceptual organizing principle}: it does not modify the formalism, derive new results, or generate predictions beyond standard quantum theory (see Section~\ref{sec:foundations}, ``Scope and Formal Demarcation''). The value lies in revealing structural parallels across scales, not in replacing established formalisms.

\subsection{States as Strategies, Superposition as Mixed Strategy}

We interpret quantum states directly as strategies in this thermodynamic game. A pure quantum state $|\psi_i\rangle$ corresponds to a pure strategy. Consequently, a quantum superposition, defined as $\sum_i c_i |\psi_i\rangle$, is identified as a mixed strategy. The squared modulus of the amplitude, $|c_i|^2$, represents the probability weight assigned to that specific strategy. This identification operates within the commutative specialization established in Section~\ref{sec:foundations}: the diagonal density matrix $\rho = \sum_i |c_i|^2 |i\rangle\langle i|$ reproduces the classical mixed-strategy structure, while off-diagonal coherences require a noncommutative extension beyond the scope of the present analysis.

In the mathematical formalism of stochastic game dynamics, a mixed strategy allows an agent to navigate uncertainty by maintaining a probability distribution over the available action space \citep{BusemeyerBruza2012}. The ability to maintain superpositions allows the field excitations to simultaneously probe multiple evolutionary pathways, ensuring that the entire configuration space is mapped before a thermodynamic commitment is mandated by environmental interaction.

Crucially, superposition does not represent epistemic uncertainty regarding a hidden underlying reality, but rather the objective strategic coexistence of multiple field configurations prior to payoff realization.

\subsection{The Born Rule as Equilibrium Distribution}

Within this game-theoretic construct, the Born rule is not an independent postulate. Instead, the Born rule can be reinterpreted as the equilibrium distribution of strategies in a mixed-strategy Nash equilibrium (see Clarification below for the precise scope of this claim). The payoff for the system is the local entropy production at the moment of measurement (interaction). Strategies with a higher expected thermodynamic payoff are realized with greater frequency. Therefore, the probability distribution defined by $|\psi|^2$ is the stationary distribution of the game.

\subsubsection{Reinterpretation via Envariance}

\textbf{Clarification:} The following does not constitute an independent derivation of the Born rule. It is a \emph{reinterpretation} of Zurek's envariance program \citep{Zurek2003,Zurek2005} within the game-theoretic vocabulary of FST-I. Zurek's original derivation is itself contested: Schlosshauer \& Fine (2005) \citep{SchlossFine2005} raise objections to the claim that the Born rule follows from envariance alone without circularity. The contribution of FST-I here is to observe that Zurek's symmetry arguments map naturally onto the Nash equilibrium concept---a conceptual connection, not a new mathematical result.

The derivation of the Born rule as an equilibrium distribution is illustrated through the mechanism of environment-assisted invariance, or ``envariance,'' pioneered by Zurek \citep{Zurek2003,Zurek2005}.

When a quantum system $\mathcal{S}$ becomes entangled with its surrounding environment $\mathcal{E}$, the composite pure state can be expressed via the Schmidt decomposition:
\begin{equation}
|\Psi_{SE}\rangle = \sum_k \alpha_k |s_k\rangle \otimes |e_k\rangle
\end{equation}
where $\{|s_k\rangle\}$ and $\{|e_k\rangle\}$ are orthonormal bases for the respective Hilbert spaces $\mathcal{H}_S$ and $\mathcal{H}_E$ \citep{SchlossFine2005}.

Envariance dictates that a local property of the system is independent of transformations that can be perfectly undone by acting solely on the environment:

\begin{center}
\begin{tabularx}{\textwidth}{@{}>{\raggedright\arraybackslash}p{3.3cm}>{\raggedright\arraybackslash}p{4.1cm}Y@{}}
\toprule
\textbf{Transformation} & \textbf{Mechanism} & \textbf{Implication} \\
\midrule
Phase Envariance & $u_S = e^{i\beta_k}|s_k\rangle\langle s_k|$ & Phases carry no strategic information \\
Swap Envariance & Swapping $|s_1\rangle$ and $|s_2\rangle$ & Equal amplitudes $\Rightarrow$ equal probabilities \\
Counting Extension & Fine-graining environment & Generalizes to all rational $p_k$ \\
\bottomrule
\end{tabularx}
\end{center}

The frequency of strategy realization is deterministically proven to be $p_k \propto |\alpha_k|^2$. In FST-I, this envariant derivation is recognized as the formal signature of a Nash equilibrium. Because the environment continuously monitors the system, any strategic deviation from the Born probability distribution would be sub-optimal and immediately corrected by the rapid suppression of phase coherences. A related relaxation mechanism---though based on fundamentally different dynamics---appears in the pilot-wave interpretation: Valentini \& Westman \citep{ValentiniWestman2005} show that non-Born distributions relax to $|\psi|^2$ through mixing by the pilot wave in de~Broglie--Bohm theory. Their result supports the universality of Born-rule relaxation across interpretations, but the underlying mechanism (deterministic hidden-variable dynamics) differs from the game-theoretic picture proposed here.

\subsection{The Path Integral as a Sum Over Strategies}

The conceptual anchor of this interpretation is the Feynman path integral formulation. The transition amplitude is given by the sum over all possible histories:
\begin{equation}
A = \sum_{\text{paths}} e^{\frac{i}{\hbar} S[\text{path}]}
\end{equation}

Here, each individual path represents a complete strategy, and the action $S$ encodes the cumulative payoff of that strategy over time. Quantum interference acts as the mechanism of strategic competition. In the semiclassical limit, only paths near the stationary phase contribute constructively.

\paragraph{Stationary phase as saddle-point equilibrium.}
An important subtlety must be addressed: in Minkowski signature, the stationary-phase condition $\delta S = 0$ yields a \emph{saddle point}, not a maximum. The correct game-theoretic analogue is therefore not single-agent payoff maximization, but a \emph{zero-sum saddle-point equilibrium} (minimax Nash). In first-order (Hamiltonian) form,
\begin{equation}
S[q,p] = \int_{t_0}^{t_1} \left( p \cdot \dot{q} - H(q,p) \right) dt,
\end{equation}
the classical equations are equivalent to the saddle-point conditions $\delta S / \delta q = 0$, $\delta S / \delta p = 0$---no infinitesimal unilateral variation in either conjugate variable can improve the corresponding side of the saddle.

\paragraph{Euclidean formulation.}
After Wick rotation ($t \mapsto -i\tau$), the weight becomes $\exp(-S_E[q]/\hbar)$, and the semiclassical limit is governed by \emph{minima} of $S_E$. A literal payoff interpretation is then obtained by defining the utility as $U[q] := -S_E[q]$, so that equilibrium paths maximize $U$ (equivalently minimize $S_E$). The Minkowski stationary-phase statement is recovered by analytic continuation. In the remainder of this paper, references to ``payoff maximization'' should be understood in this Euclidean sense unless stated otherwise.

This formulation effectively replaces the classical notion of a single, unique trajectory with a functional integral over an infinity of quantum-mechanically possible trajectories \citep{PathIntegral}. When applied to field excitations, the Euclidean action serves as the utility function \citep{PathIntegralOptimization}. The entire quantum propagation process is a parallelized mathematical computation where every conceivable strategy is evaluated simultaneously.

\subsection{Destructive Interference as Strategy Elimination}

Non-stationary paths---those far from the saddle point of the action---experience rapid phase oscillation, leading to destructive interference. In the context of game theory, destructive interference plays the exact role of evolutionary elimination: strategies that deviate from the equilibrium cancel each other out. What remains are the stable, constructive coalitions of paths that represent the robust equilibria of the field dynamics.

This dynamic is loosely analogous to selection processes in evolutionary biology, insofar as non-optimal strategies are eliminated. However, the analogy should not be over-read: destructive interference operates on complex amplitudes in a single system, not on populations with inheritance, variation, and differential fitness. In a highly competitive stochastic environment, strategies that deviate from the saddle-point equilibrium exhibit wildly divergent complex phases. When summed in the path integral, these highly oscillating terms sum to zero, effectively erasing the probability of the system occupying non-equilibrium states.

\subsection{Measurement as Payoff Realization}

The controversial concept of wave function collapse is reframed as payoff realization. Measurement corresponds to payoff realization through irreversible coupling to macroscopic, thermodynamic degrees of freedom. There is no mystical ``physical collapse''; rather, the measurement event marks the end of strategic flexibility. The system commits to a definitive state, extracting the thermodynamic payoff.

Decoherence theory provides the exact physical mechanism for this payoff realization \citep{MeasurementClassical}. As a quantum system interacts with a macroscopic environment, the phase relations of the superposition become irreversibly delocalized, transferring information into the unobservable degrees of freedom of the surrounding bath. This process, einselection, imposes an effective ban on the vast majority of the Hilbert space, eliminating fragile ``Schr\"odinger cat'' states and leaving only a discrete set of robust, classical pointer states \citep{Zurek2003}.

Furthermore, the Heisenberg uncertainty principle is understood as a requirement for strategic flexibility: conjugate variables cannot be simultaneously fixed because doing so would over-constrain the strategy space, preventing the system from reaching an optimal mixed-strategy equilibrium \citep{HeisenbergCovariant}.

\subsection{Scope and Formal Demarcation}

It is vital to state explicitly what this interpretation entails and what it avoids. This framework does not modify the mathematical formalism of quantum mechanics. It does not introduce hidden variables, nor does it alter the predictive outcomes of standard quantum theory. Instead, it provides a functional interpretation in which the axiomatic structure of quantum mechanics follows necessarily from optimization and equilibrium principles governing thermodynamic systems.

\paragraph{Value of the game-theoretic reformulation.}
We do not claim a new derivation of the Born rule or new quantitative predictions beyond standard quantum theory. The contribution is a \emph{structural unification}: envariance selects Born weights as a stability/fixed-point condition under admissible transformations, in direct analogy to KL/free-energy equilibria in statistical mechanics (Gibbs minimum, FST-II) and replicator dynamics in evolutionary biology (Nash/ESS, FST-III). The game-theoretic vocabulary makes the implicit assumptions of the envariance route explicit as ``rules of interaction''---allowed symmetries, information access, and stability under unilateral deviations---thereby improving axiom transparency. In this sense the contribution is conceptual unification and axiom bookkeeping, not a new theorem.

By anchoring the postulates of quantum mechanics directly to the formalisms of stochastic game dynamics and envariance, FST-I provides a common language that connects quantum measurement, chemical selection, and biological evolution as instances of a shared fixed-point/variational template \citep{SchlossFine2005}. The mathematical apparatus of quantum mechanics is preserved in its entirety \citep{QuantumOptimization}, but its ontological status is shifted: it is recognized as an optimal control framework for thermodynamic efficiency under conditions of incomplete information.

% ===================================================================
\section{The Maximum Entropy Production Principle as Global Utility Function}
\label{sec:mepp}
% ===================================================================

To transition from the micro-dynamics of quantum mixed strategies to the macro-dynamics of cosmological evolution, we must explicitly define the global utility function driving the thermodynamic game. In FST-I, this function is provided by the Maximum Entropy Production Principle (MEPP) \citep{MEPP2004}.

\subsection{Statement of MEPP}

MEPP asserts that complex, non-equilibrium systems subjected to steady external forcing will inevitably self-organize to select the specific thermodynamic path that maximizes the rate of entropy production ($\dot{S}$) \citep{MEPPFoundations}. Originally formulated in the context of planetary climatology by Paltridge \citep{Paltridge1975} and subsequently generalized across fluid dynamics, chemistry, and biology \citep{MEPP2004}, MEPP represents the continuous drive of the universe to dissipate energy gradients as efficiently as possible.

\subsection{Status and Controversy of MEPP}

\textbf{Important:} MEPP is an active subject of debate in the non-equilibrium thermodynamics community. Its status as a fundamental law, as a useful heuristic, or as a framework-dependent postulate is contested:

\begin{itemize}
    \item \textbf{Grinstein \& Linsker (2007)} demonstrated that MEPP does not hold in general for discrete stochastic systems and cannot be derived from standard non-equilibrium thermodynamics without additional assumptions \citep{GrinsteinLinsker2007}. Explicit counterexamples exist: systems governed by Ginzburg--Landau-type dynamics can select states of \emph{minimal} rather than maximal entropy production.
    \item \textbf{Polettini (2013)} sharpened this objection for Ziegler-style maximum entropy production beyond linear response: stochastic-thermodynamic models can violate the reciprocal relations implied by thermodynamic orthogonality, and network steady states are not produced by entropy-production maximization alone~\citep{Polettini2013}.
    \item \textbf{Dewar (2003)} argued for MEPP as a consequence of maximum entropy inference (Jaynes' maximum entropy formalism \citep{Jaynes1957} applied to path probabilities), but this derivation has been criticized on technical grounds \citep{Dewar2003}.
    \item \textbf{Martyushev \& Seleznev (2006)} provide the most comprehensive review, concluding that MEPP has substantial empirical support across many systems but lacks a universally valid derivation \citep{MartyushevSeleznev2006}.
    \item \textbf{Sawada, Daigaku \& Toma (2025)} recently reviewed MEPP in the context of the origin and evolution of life, confirming that MEPP describes macroscopic evolutionary trends well, but a rigorous first-principles derivation far from equilibrium remains incomplete~\citep{Sawada2025}.
\end{itemize}

In FST-I, MEPP is treated as a working hypothesis---a plausible organizing principle with empirical support in multiple domains---rather than a proven axiom. The core claim of the paper does not require MEPP to hold universally; it requires that the Standard Model parameter set is \emph{consistent with} a constrained entropy-production optimum. Whether this is the result of MEPP ``selecting'' the parameters over cosmological timescales or is a structural property of stable physics, the falsifiability criteria in Section~\ref{sec:predictions} are designed to adjudicate. More recently, Martyushev (2010) \citep{Martyushev2010} identified two basic questions that any MEPP application must answer: (i)~whether the system satisfies local thermodynamic equilibrium, and (ii)~whether the entropy production is bilinear in fluxes and forces. For the far-from-equilibrium systems considered in FST-I, the second condition is generally not met, which restricts MEPP to a heuristic stability filter rather than a rigorous selection principle.

\paragraph{Circularity caveat.} When MEPP is used as the payoff function and Nash equilibria are defined as MEPP-maximizing configurations, the claim that observed systems are Nash equilibria because they maximize entropy production becomes definitionally circular. The non-circular content of the framework lies in the \emph{stability} claim: that the observed parameters are not merely entropy-producing but second-order stable against perturbations. This stability property is testable independently of the payoff definition (Section~\ref{sec:finetuning}).

\subsection{Operational Modes}

In multi-stable stochastic systems, MEPP manifests in two distinct operational modes \citep{MartyushevSeleznev2006}:
\begin{enumerate}
    \item \textbf{State selection principle:} Determines which specific steady state is occupied in a highly degenerate landscape;
    \item \textbf{Gradient response principle:} Dictates how the system reacts to perturbations.
\end{enumerate}

When mapped onto fundamental physics, MEPP requires that the rules of the game---the fundamental constants of nature---must themselves be structured to permit the highest possible sustained thermodynamic dissipation \citep{EntropyMaximum}. The field excitations act as agents whose interaction cross-sections and coupling strengths are constrained by this thermodynamic utility function. In the FST-I framework, this does not imply that physical constants vary over time; rather, the constants are selected (or filtered) at early cosmological epochs and remain fixed thereafter.

\begin{remark}[MEPP as Stability Filter, not Maximization Principle]
\label{rem:mepp-stability-filter}
A crucial refinement emerges from the second-order resolvent analysis developed in the gauge-theoretic reformulation of the Riemann Hypothesis \citep{GeigerRH2026b}. The FST framework describes \emph{stability principles}, not maximization principles. Leading-order dynamics are degenerate: many configurations of the Standard Model parameters yield comparable entropy production rates at first order. Selection occurs at second order via resolvent-type coupling between degenerate and complementary subspaces. MEPP thus acts as a \emph{stability filter} that identifies the unique resolvent-stable fixed point within the degenerate manifold, rather than selecting a global maximum. The observed Standard Model parameters are not global entropy-production maxima, but second-order-stabilized equilibria---configurations that are robust against perturbations propagating through resonance couplings between parameter sectors.
\end{remark}

\paragraph{Scale separation and the FEP--MEPP duality.}
A related concern is the apparent tension between entropy \emph{minimization} (Friston's free energy principle, FEP) and entropy \emph{maximization} (MEPP). In the FST framework this tension dissolves through scale separation: each parameter sector locally minimizes its internal variational free energy, thereby maintaining structural integrity against fluctuations---the field-theoretic analogue of a Markov blanket \citep{Friston2010}. These individually stable sectors, however, collectively constitute a highly efficient dissipative architecture whose \emph{aggregate} effect is maximal entropy production at the macroscopic scale. Internal FEP minimization and external MEPP maximization are therefore not competing principles but complementary aspects of the same multi-scale equilibrium; the formal treatment is developed in FST-III (Proposition~\texttt{fep-mepp}).

\begin{conjecture}[FEP--England Duality]\label{lem:fep-england}
Let $(I, B, E)$ be a Markov blanket partition of a system at nonequilibrium steady state satisfying local detailed balance. Let $\mathcal{A}_E$ denote the set of admissible blanket-state trajectories over $[0, T]$, and let $\mathcal{Q}_I$ denote the set of variational densities over internal states. Then, in expectation over trajectories,
\[
  \max_{a \in \mathcal{A}_E} \langle \Delta S_{\mathrm{ext}} \rangle \;\Longleftrightarrow\; \min_{q \in \mathcal{Q}_I} F_{\mathrm{int}},
\]
where $\Delta S_{\mathrm{ext}} \geq 0$ is the entropy produced in the external states, bounded below by the Crooks fluctuation theorem \citep{Crooks1999}, and $F_{\mathrm{int}} = \mathbb{E}_q[\ln q(s) - \ln p(s,o)]$ is the variational free energy of the internal states \citep{FristonFreeEnergy}. The two optimization problems operate on \emph{different} variable sets ($\mathcal{A}_E$ vs.\ $\mathcal{Q}_I$); the claimed equivalence asserts that at the steady-state fixed point, the solution of one entails the solution of the other.
\end{conjecture}

This duality is stated without proof; the argument is a heuristic motivated by the Crooks fluctuation theorem, England's dissipative self-replication bound \citep{England2013}, and the variational formulation of the Free Energy Principle. A rigorous derivation under the stated conditions has not been published, and the precise sense of the biconditional (whether it holds for all steady states or only for a restricted class) remains an open question. England's original result provides a \emph{lower bound} on entropy production for self-replicators, not an equivalence with free-energy minimization; the conjectured duality is therefore stronger than what the cited literature supports.

\begin{remark}
This duality holds in expectation over trajectories, not pathwise. It requires local detailed balance at the blanket boundary and does not extend to systems with broken time-reversal symmetry. The duality should therefore be understood as a statistical organizing principle rather than a dynamical law.
\end{remark}

% ===================================================================
\section{Fine-Tuning as Optimization}
\label{sec:finetuning}
% ===================================================================

The proposition that the Standard Model is optimized for entropy production must be operationalized to yield testable physical consequences. Historically, the exact values of the fundamental constants have been viewed as an anthropic coincidence, fine-tuned to an arbitrary degree to permit the existence of observers \citep{Adams2019}. FST-I rejects this passive interpretation. The parameter landscape is not tuned for life; it is tuned for optimal, maximal thermodynamic throughput, of which biological complexity is merely a localized, highly dissipative byproduct.

\subsection{Definition of the Optimization Functional}

To formulate the Maximum Entropy Production Principle analytically, we define the total entropy production functional over the relevant cosmological epoch:
\begin{equation}
\mathcal{S}[\vec{p}] = \int_{t_{BB}}^{t_\Lambda} \dot{S}(\vec{p}, t) \, dt
\end{equation}

Here, $\vec{p}$ denotes the set of dimensionless Standard Model parameters, such as $\vec{p} = \{\alpha_{EM}, \alpha_s, \alpha_W, m_e/m_p, m_u/m_d, \ldots\}$. The integration boundaries span from the Big Bang ($t_{BB}$) to the onset of dark energy domination ($t_\Lambda$), defining the primary epoch of structure formation and stellar nucleosynthesis.

Crucially, $\dot{S}$ does not represent the coarse-grained total entropy of the universe (which is dominated by the Bekenstein-Hawking entropy of supermassive black holes, $\sim 10^{104}\,k_B$ \citep{EntropyMaximum}). Black holes are thermodynamic endpoints, not active production mechanisms. Instead, $\dot{S}$ is defined as the baryon-weighted entropy production rate generated by sustained nuclear, electromagnetic, and gravitational processes---primarily driven by stellar fusion and subsequent photon thermalization.

\subsection{Variational Principle with Constraints}

If the Standard Model is the result of a thermodynamic game, its parameter configuration must satisfy a variational principle $\delta \mathcal{S} = 0$ with $\delta^2 \mathcal{S} \leq 0$ (constrained local maximum), subject to specific boundary conditions:
\begin{enumerate}
    \item \textbf{Nuclear Stability:} The existence of bound, stable nuclei (requiring specific ratios of strong to electromagnetic couplings).
    \item \textbf{Proton Stability:} The proton must be stable against decay within the Standard Model (baryon number conservation), ensuring that hydrogen fuel persists over cosmological timescales ($\tau_p \gg t_\star$).
    \item \textbf{Chemical Complexity:} The possibility of stable electron orbitals to form atomic and molecular networks.
\end{enumerate}

\subsection{The Diproton Catastrophe}

Detailed analyses reveal extreme parameter sensitivity. As demonstrated by Adams (2019), the strong coupling constant ($\alpha_s$) is subject to highly sensitive fine-tuning \citep{Adams2019}:

\begin{itemize}
    \item If the strong force were \textbf{reduced by $\sim$10\%}, the deuterium nucleus would become unbound, completely severing the primary nucleosynthetic pathway for stellar fusion.
    \item If the strong force were \textbf{increased by $\sim$10\%}, diprotons (helium-2) would become bound and highly stable. Stars would bypass the sluggish weak interaction entirely and burn through their hydrogen fuel via the strong interaction, causing main-sequence stars to explode or burn out in a fraction of their current lifespans.
\end{itemize}

By avoiding the diproton catastrophe, the universe enforces a slow, metered release of energy, ensuring that the total integrated functional $\mathcal{S}$ over the epoch $[t_{BB}, t_\Lambda]$ is globally maximized.

\subsection{The Triple-Alpha Resonance}

The production of critical heavy elements relies entirely on the triple-alpha reaction \citep{Adams2019}. The viability of this process requires the precise existence of a $0^+$ resonance energy state in the carbon-12 nucleus. If $\alpha_{EM}$ varied such that this resonance energy shifted outside a narrow window of $\sim 100$ keV to $\sim 400$ keV, stellar nucleosynthesis would fail to produce the carbon and oxygen necessary for complex chemistry and dust formation, radically altering the opacity and cooling mechanisms of galaxies.

\subsection{Concrete Toy Variation: The Electron Mass}

Consider a heuristic variation of the electron mass, $m_e \to m_e(1+\epsilon)$ with $|\epsilon| \ll 1$, while holding other parameters constant.

\textbf{Positive Variation ($m_e \uparrow$):} An increase in electron mass tightens atomic orbitals and lowers the Thomson opacity ($\kappa \propto m_e^{-2}$), making the stellar interior more transparent and raising the net luminosity. This accelerates the effective fusion rate ($L_\star \sim \Gamma_{fusion}$), leading to drastically shorter stellar lifetimes ($t_\star$). While the instantaneous rate $\dot{S}$ might spike, the integrated functional $\mathcal{S}[\vec{p}]$ over the cosmological timeframe drops significantly due to premature stellar burnout.

\textbf{Negative Variation ($m_e \downarrow$):} A decrease in electron mass increases atomic radii and \emph{increases} Thomson opacity ($\kappa \propto m_e^{-2}$, so lower $m_e$ yields higher $\kappa$), making stellar envelopes more opaque. While the Gamow energy $E_G = (\pi \alpha)^2 m_r c^2$ depends on the reduced nuclear mass and not directly on $m_e$, the causal chain operates indirectly: higher opacity traps radiation more effectively, which raises the equilibrium core temperature $T_c$, which in turn increases the Gamow tunneling rate $\Gamma_{\text{fusion}} \propto \exp(-\sqrt{E_G / k_B T_c})$. However, the increased opacity simultaneously reduces the photon diffusion rate and thereby lowers the net luminosity. The resulting longer main-sequence lifetimes lead to higher integrated entropy production---consistent with the monotonic increase of $\mathcal{S}_{\text{total}}$ with decreasing $m_e$ found in Section~\ref{sec:stellar_model}. The observed $m_e$ is therefore not fixed by stellar entropy production alone but by multi-scale constraints from atomic and chemical physics (see Section~\ref{sec:stellar_model} for quantitative details).

In summary, the observed fine-structure constant lies near the optimum of $\mathcal{S}$ for fixed $m_e/m_p$ (Table~\ref{tab:alpha_scan}: 98.8\% of maximum), whereas $m_e/m_p$ is constrained by multi-scale physics rather than by a single-scale stellar entropy peak---a dichotomy quantified in Section~\ref{sec:stellar_model}.

\subsection{Connection to the Nash Equilibrium}

This variational extremum maps directly onto a Nash equilibrium. Each fundamental parameter in $\vec{p}$ maps onto a ``strategy'' in the formal game-theoretic sense---i.e., a configuration variable whose value is selected by the stability condition.

Because the parameters are highly coupled (e.g., stellar fusion requires the interplay of the strong force for binding, the weak force for proton-to-neutron conversion, and electromagnetism for radiation), the extremum of $\mathcal{S}$ is a cooperative state. The non-trivial content of this Nash characterization lies not in the payoff definition (which would be circular; see the circularity caveat in Section~\ref{sec:mepp}) but in the \emph{stability} claim: any unilateral variation of a single parameter triggers a disproportionate collapse in the coupled structural prerequisites (nuclear stability, atomic binding, stellar longevity), a property that is independently verifiable through parameter-sensitivity studies \citep{Adams2019}.

This framework replaces static, anthropic fine-tuning with dynamic, cooperative stability.

\begin{assumption}[Potential Game Structure]\label{ass:potential-game}
The thermodynamic game defined by the sector payoffs $u_i(p) = \dot{S}_i(p)$ admits a \emph{potential function} $S(p)$ such that
\[
  u_i(p) = \frac{\partial S}{\partial p_i}(p) \quad \forall\, i.
\]
Under this assumption, best-response dynamics reduce to gradient ascent on~$S$, the Hessian $H_{ij} = \partial_{p_i}\partial_{p_j} S$ serves as the Jacobian of the dynamics, and resolvent stability is well-defined. Without this assumption, $H$ is not a physically meaningful stability operator.
\end{assumption}

\begin{remark}[Physical Motivation for the Potential Structure]
\label{rem:potential-motivation}
Near thermodynamic equilibrium, the Onsager reciprocal relations guarantee that the matrix of transport coefficients is symmetric, $L_{ij} = L_{ji}$, which implies that the entropy production $\dot{S} = \sum_{ij} L_{ij} X_i X_j$ derives from a quadratic potential in the thermodynamic forces $X_i$ \citep{Prigogine1967}. The potential game assumption extends this structure to the far-from-equilibrium regime, where Onsager reciprocity no longer holds in general and the existence of a global potential function is not guaranteed \citep{GrinsteinLinsker2007}. This extension is the central untested structural assumption of the framework. Empirical support comes from the observation that in known physical systems with high symmetry (e.g., the Standard Model gauge structure), the coupling matrix between sectors exhibits approximate symmetries that make a potential decomposition plausible---but a rigorous derivation for the full Standard Model parameter space remains an open problem.
\end{remark}

\begin{remark}[Fine-Tuning as Resolvent Stability]
\label{rem:fine-tuning-resolvent}
The resolvent perspective \citep{GeigerRH2026b} provides a structural explanation for fine-tuning that is neither teleological nor anthropic. At leading order, small parameter variations are neutral---many nearby configurations produce comparable entropy. However, these variations destabilize the system at second order through resonance couplings between parameter sectors (analogous to the second-order spectral splitting observed in the Weil quadratic form analysis, where the stability gap emerges not at leading order but through coupling between complementary subspaces). Fine-tuning is thus the signature of resolvent stability: the observed parameters occupy the unique fixed point where all second-order cross-couplings between degenerate modes are simultaneously stabilized. The ``fine-tuned'' appearance is an artifact of projecting a multi-dimensional stability condition onto one-parameter slices.
\end{remark}

\subsection{Scope and Limitations}

We do not claim a full derivation of the Standard Model from first principles or an a priori calculation of the exact value of $\alpha_{EM}$. Rather, we claim that the parameter values correspond to a constrained, entropy-maximizing equilibrium.

\textbf{Critical limitation:} While Section~\ref{sec:stellar_model} provides a first semi-analytic computation demonstrating that $\mathcal{S}_{\text{total}}(\alpha_{\text{obs}})$ achieves 98.8\% of the unconstrained stellar maximum, this calculation addresses only a single entropy channel (main-sequence stellar radiation) and varies only two parameters ($\alpha$ and $m_e/m_p$). A complete test of the FST-I hypothesis requires evaluating $\mathcal{S}[\vec{p}] = \int_{t_{BB}}^{t_\Lambda} \dot{S}(\vec{p}, t)\,dt$ across all 19 Standard Model parameters and all entropy production channels (stellar, gravitational, chemical, nuclear). Such a computation would involve: (i) stellar evolution codes with non-standard nuclear parameters, (ii) galaxy formation simulations, and (iii) integration over the full cosmological history.

The semi-analytic result is encouraging---the observed $\alpha$ is near-optimal for stellar entropy---but does not yet constitute a full demonstration. In particular, the monotonic $m_e/m_p$ dependence shows that multi-scale constraints (atomic stability, chemistry viability) must be incorporated. Extending this calculation is the primary task for future work.

FST-I is a functional framework. It does not dictate the fundamental microscopic topology of any deeper theory. Instead, it places an absolute boundary condition: the effective low-energy parameters that emerge from any fundamental theory must satisfy the MEPP variational principle to remain stable over cosmological timescales.

% ===================================================================
\section{Symmetry Breaking as Phase Transition in the Game}
\label{sec:symmetry}
% ===================================================================

Section~\ref{sec:quantum} established quantum mechanics as the equilibrium distribution of strategies, and Section~\ref{sec:finetuning} defined the optimization functional for the Standard Model parameters. This section addresses the dynamic nexus between them: the origin of mass and structure.

\subsection{Symmetry as a Degenerate Equilibrium}

In the game-theoretic framework, a symmetric phase corresponds to a degenerate Nash equilibrium: multiple strategies yield identical thermodynamic payoffs. The unbroken symmetry represents a flat payoff plateau in the configuration space. Because no single strategy is thermodynamically preferred, the system is marginally stable.

During the earliest epochs of the universe, prior to the electroweak epoch, the ambient thermal energy density vastly exceeded the rest mass energies of all interacting particles \citep{HiggsVEV}. In this high-temperature regime, maintaining a massless, highly symmetric plasma is the optimally efficient configuration for rapid energy dissipation, as it maximizes the available degrees of freedom and collision rates \citep{HiggsVEV2024}.

\subsection{Electroweak Symmetry Breaking as Instability}

As the external cosmological constraints evolve---through the expansion and cooling of the universe---the previously symmetric equilibrium becomes unstable. Above the critical temperature ($T > T_c$), the minimum of the effective potential sits at $\phi = 0$, corresponding to the unique stable equilibrium. As the temperature drops below $T_c$, $\phi = 0$ transitions from a stable minimum to a local maximum (unstable equilibrium), and new minima appear at $\phi \neq 0$.

The self-interacting scalar field potential that governs this transition is:\footnote{We follow the Peskin \& Schroeder convention \citep{HiggsVEV} where $V(\phi) = -\mu^2 |\phi|^2 + \lambda |\phi|^4$ with $\mu^2 > 0$ for the broken phase. In the real-scalar shorthand used here, the factor $\frac{1}{2}$ and $\frac{1}{4}$ are conventional normalizations.}
\begin{equation}
V(\phi) = -\frac{1}{2}\mu^2 \phi^2 + \frac{1}{4}\lambda \phi^4
\end{equation}

As the universe cools, the mass parameter $\mu^2$ evolves, passing through zero and becoming positive (in the convention where the negative sign is explicit). This provokes spontaneous symmetry breaking \citep{HiggsVEV2024}. In the context of the thermodynamic game, the vacuum's transition to a broken-symmetry phase is an inevitable dynamical response driven by the MEPP functional.

\subsection{The Higgs VEV as an Attractor of the Game}

The Higgs vacuum expectation value (VEV) corresponds to the stable fixed point---the ultimate attractor---of the Level-1 game. The VEV is not merely ``spontaneously chosen'' by a random fluctuation; it is dynamically selected because a non-zero expectation value, $\langle\phi\rangle \neq 0$, strictly maximizes the local payoff.

The exact numerical value of the Higgs VEV, approximately 246 GeV, serves as the fundamental mass generator for the W and Z bosons, as well as the fermions via Yukawa couplings \citep{HiggsVEV}. This precise scaling parameter separates the massless photon from the massive mediators of the weak force, effectively creating the disparity in force strengths that enables long-term structural stability of nuclear matter.

Masses emerge specifically because they enable higher-order, sustained entropy production. A universe populated exclusively by massless particles would remain radiation-dominated indefinitely, preventing gravitational collapse into bound structures (stars, galaxies) and rendering it thermodynamically inefficient for sustained dissipation.

\subsection{Connection to Cosmological Saturation}

This dynamic shares a deep formal isomorphism with the dynamics at the cosmological scale. Recent cosmological reconstructions linking the kinematic jerk parameter ($j = 1$ for $\Lambda$CDM) to scalar field dynamics have demonstrated a theoretical cosmic variation in the proton-to-electron mass ratio ($\mu = m_p/m_e$; note that the rest of this paper uses the reciprocal $m_e/m_p$) driven by the evolution of the Higgs VEV over time ($v(z)$) \citep{HiggsVEV2024}. Theoretical calculations of this variation fit within observational constraints from quasar absorption spectra, providing evidence that the microscopic symmetry-breaking order parameter ($v$) remains dynamically coupled to the macroscopic expansion rate.

In the Curvature Relaxation Model (CRM \cite{GeigerCRM2026}, see also \cite{FST-Unified2026}), the coarsest-scale (``Level-0'') spacetime alignment yields a $\tanh$ saturation curve as its mean-field solution. At the next finer scale (``Level~1''), the quartic Higgs potential serves as the effective description. Both phenomena arise from the same underlying optimization logic operating under different scales of coarse-graining.

\subsection{Cosmological Phase Transitions as Successive Nash Equilibria}

Each cosmological phase transition corresponds to a transition between Nash equilibria:

\begin{center}
\begin{tabularx}{\textwidth}{@{}>{\raggedright\arraybackslash}p{3.1cm}YY@{}}
\toprule
\textbf{Transition} & \textbf{Strategic Shift} & \textbf{Consequence} \\
\midrule
Electroweak & Settles into VEV configuration & Enables gravitational clustering \\
QCD Confinement & Settles into bound-state configuration & Creates stable proton repositories \\
Recombination & Settles into neutral-atom configuration & Creates CMB thermal sink \\
\bottomrule
\end{tabularx}
\end{center}

Every phase transition represents a reduction in pure strategic freedom (lowering of symmetry) in exchange for an increase in irreversible, robust entropy production.

% ===================================================================
\section{The Role of Information and Decoherence}
\label{sec:decoherence}
% ===================================================================

While Section~\ref{sec:quantum} established the mathematical equivalence of quantum superposition to mixed strategies, it is necessary to examine the precise informational mechanics driving strategy selection. In standard quantum theory, the environment acts as a passive reservoir; in FST-I, the environment acts as an active, competitive opponent in a zero-sum informational game \citep{Zurek2003}.

\subsection{Decoherence as Strategic Enforcement}

Decoherence theory provides the exact dynamic mechanism by which this game is resolved. As a quantum system evolves, it continuously interacts with the surrounding vacuum and thermal baths. This interaction entangles the system's strategic state with the unobservable degrees of freedom of the environment \citep{MeasurementClassical}.

Because the local observer is barred from accessing the entire universal wave function, the reduced density matrix of the system rapidly loses its off-diagonal interference terms through tracing over environmental states \citep{Lindblad1976}. For mesoscopic systems, decoherence timescales range from ${\sim}10^{-23}\,$s (strongly interacting hadrons in a thermal bath) to ${\sim}10^{-17}\,$s (molecular superpositions at room temperature), ensuring that strategy selection occurs on timescales far shorter than any macroscopic observation.

\subsection{Einselection as Nash Enforcement}

This process, termed einselection (environment-induced superselection), is not a passive decay; it is the systematic enforcement of the Born rule equilibrium \citep{Zurek2003}. The environment selectively monitors certain observables (typically those commuting with the interaction Hamiltonian, such as position), suppressing any mixed strategies (superpositions) that do not align with stable ``pointer states.''

The realization of a specific measurement outcome is therefore the thermodynamic culmination of the game, wherein the informatic potential of a vast, uncollapsed mixed strategy is irreversibly converted into a singular physical state, generating a localized spike in entropy and securing the thermodynamic payoff \citep{Lindblad1976}.

% ===================================================================
\section{Connections to Existing Approaches}
\label{sec:connections}
% ===================================================================

Having established the core architecture of FST-I---quantum mechanics as mixed-strategy equilibrium (Section~\ref{sec:quantum}), MEPP as the global payoff function (Section~\ref{sec:mepp}), fine-tuning as constrained optimization (Section~\ref{sec:finetuning}), symmetry breaking as phase transition (Section~\ref{sec:symmetry}), and decoherence as strategic enforcement (Section~\ref{sec:decoherence})---we now situate this framework within the broader landscape of related approaches.

\subsection{Entropic Gravity: Verlinde and Padmanabhan}

Verlinde's entropic gravity proposal \citep{Verlinde2011} suggests that gravitational forces arise from entropic gradients. Padmanabhan's program \citep{Padmanabhan2010} derives gravitational field equations from thermodynamic identities. Our framework extends this logic to the full Standard Model parameter space.

\subsection{Environment-Induced Superselection: Zurek}

Zurek's decoherence program \citep{Zurek2003} explains classical emergence through environmental monitoring. Our framework reinterprets pointer states as Nash-stable strategies, with einselection as the mechanism by which non-optimal strategies are eliminated.

\subsection{Free Energy Principle: Friston}

Friston's free energy principle \citep{Friston2010} proposes that biological systems minimize variational free energy. We view this as a special case operating at the biological scale, inheriting its structure from particle-level optimization through coarse-graining.

\paragraph{Active Inference as Stochastic Optimal Control.}
Active Inference can be embedded in the MFG framework as a stochastic optimal control problem \citep{LasryLions2006,LasryLions2007}. Each agent maintains internal beliefs $\mu_t \in \mathbb{R}^n$ (sufficient statistics of $q(s|\mu)$) and selects actions $a_t \in \mathcal{U}$ that modify blanket states. The dynamics follow
\[
  d\mu_t = f(\mu_t, a_t, m_t)\,dt + \sigma\,dW_t,
\]
where $m_t$ is the mean-field distribution over other agents' states. The cost functional is the variational free energy:
\[
  L(\mu, a, m) = \mathrm{KL}\bigl(q(s|\mu) \,\|\, p(s|o,a,m)\bigr) - \ln p(o|a,m).
\]
The Hamilton--Jacobi--Bellman equation for the value function $V(\mu,t) := \inf_{a} \mathbb{E}\bigl[\int_0^T L\,dt + \Phi(\mu_T)\bigr]$ reads
\[
  -\partial_t V = \inf_a\bigl[L(\mu,a,m) + \nabla_\mu V \cdot f(\mu,a,m)\bigr] + \tfrac{\sigma^2}{2}\Delta_\mu V.
\]
Coupled with the Fokker--Planck equation for $m$, this yields an MFG system whose fixed point $(V^*, m^*)$ corresponds to the Nash equilibrium of the multi-agent system. When the free energy decomposes additively, $F = \sum_i F_i$, the MFG reduces to a potential game (Assumption~\ref{ass:potential-game}), providing the formal bridge between Active Inference and the FST framework.

\subsection{Fine-Tuning Studies: Adams, Tegmark}

Adams \citep{Adams2019} and Tegmark et al.\ \citep{Tegmark2006} have systematically explored the parameter space of possible universes, mapping viable regions where stable nuclei, atoms, and stars can exist. Tegmark et al.\ identify 31 dimensionless parameters and demonstrate that many must lie within narrow ranges for structure to form. Our contribution is to provide the organizing principle: the Standard Model parameters lie not just within the viable region but near its thermodynamic optimum, as quantitatively supported by the alpha-scan results in Section~\ref{sec:stellar_model}.

\subsection{Alternative Frameworks and Their Relation to FST-I}

Several alternative explanations for fine-tuning must be considered and compared to FST-I:

\begin{itemize}
    \item \textbf{Minimum Entropy Production (MinEP):} Prigogine's minimum entropy production theorem \citep{Prigogine1967} applies to systems near thermodynamic equilibrium. It predicts the opposite of MEPP for near-equilibrium states. FST-I claims that the cosmological epoch of structure formation is far from equilibrium and thus in the MEPP regime; this claim requires verification.
    \item \textbf{Anthropic Principle:} Weinberg \citep{Weinberg1987} and others argue that observed parameters are selected by observer existence, not by optimization. FST-I makes the stronger claim that the parameters are at a thermodynamic optimum, not merely in the anthropic window. The key empirical distinction is whether the observed parameters are at the \emph{boundary} of the viable window (consistent with anthropic selection) or at an interior \emph{maximum} of $\mathcal{S}[\vec{p}]$ (consistent with FST-I). Testing this requires the numerical calculations described in Section~\ref{sec:predictions}.
    \item \textbf{String Landscape:} Susskind \citep{Susskind2005} proposes that the landscape of possible string vacua, combined with eternal inflation, provides a statistical explanation for fine-tuning. FST-I does not require string theory and makes a distinct prediction: among all viable parameter sets (not just anthropic ones), entropy production should peak near the observed values.
    \item \textbf{Cosmological Natural Selection:} Smolin \citep{Smolin1992} proposes that universes reproduce via black holes, with parameters selected to maximize black-hole production. While this provides an evolutionary selection mechanism, it does not derive dynamical equations or predict specific parameter values. FST-I replaces reproductive fitness with thermodynamic stability as the selection criterion and produces quantitative predictions.
    \item \textbf{Cellular Automaton Interpretation:} 't~Hooft \citep{tHooft2016} proposes that quantum mechanics emerges from an underlying deterministic cellular automaton dynamics. While FST-I does not invoke hidden variables, 't~Hooft's program shares the conviction that the axiomatic structure of quantum mechanics is not fundamental but derivable from deeper (in his case deterministic, in ours thermodynamic) principles. The key difference is that FST-I derives quantum structure from optimization and stability constraints without requiring microscopic determinism.
\end{itemize}

Table~\ref{tab:framework-comparison} provides a compact comparison of FST-I with the principal alternative frameworks addressing fine-tuning or the origin of physical law.

\begin{table}[ht]
\centering
\small
\caption{Comparison of frameworks addressing fine-tuning and the origin of physical parameters.}
\label{tab:framework-comparison}
\begin{tabularx}{\textwidth}{@{}>{\raggedright\arraybackslash}p{3.0cm}YYY@{}}
\toprule
\textbf{Framework} & \textbf{Mechanism} & \textbf{Predictions} & \textbf{Status} \\
\midrule
Anthropic Principle \citep{Weinberg1987} & Observer selection on $\Lambda$ & $\Lambda$ bound; no dynamics & No falsifiable mechanism \\[4pt]
Cosmological Natural Selection \citep{Smolin1992} & Black-hole reproduction selects parameters & Maximizes black-hole count & Unfalsified; no dynamics derived \\[4pt]
Free Energy Principle \citep{Friston2010} & Variational inference minimizes surprise & Agent-level behavior & Not applied to fundamental physics \\[4pt]
Cellular Automaton QM \citep{tHooft2016} & Deterministic substrates for QM & Discreteness signatures & No parameter selection \\[4pt]
\textbf{FST-I (this work)} & Stability-constrained entropy maximization & $\alpha$ within broad entropy plateau (98.8\%); $m_e/m_p$ requires multi-scale constraints (47\% in 2D) & Falsifiable via multi-parameter perturbation; pilot study only \\
\bottomrule
\end{tabularx}
\end{table}

FST-I combines (i)~a stability mechanism (Nash equilibrium of a thermodynamic game) with (ii)~a first pilot-study consistency check (the $\alpha$-scan plateau at 98.8\% entropy optimality for a single parameter in a simplified stellar model) and (iii)~an explicit falsification criterion (Section~\ref{sec:predictions}). We emphasize that (ii) is not a prediction but a post-hoc consistency check using known physics; the Gamow tunneling factor---which drives the 98.8\% result---is itself standard nuclear physics, so the near-optimality reflects the known fact that stellar fusion operates efficiently at the observed coupling strength. The anthropic principle and cosmological natural selection provide selection pressures but no dynamical equations; the free energy principle operates at the biological scale and has not been extended to fundamental constants; 't~Hooft's cellular automaton program addresses the origin of quantum mechanics but not the values of the Standard Model parameters.

The FST-I framework can, in principle, be distinguished from these alternatives through the numerical predictions in Section~\ref{sec:predictions}. However, this distinction has not yet been demonstrated computationally.

\begin{remark}[Nash Equilibria as Second-Order Stable Fixed Points]
\label{rem:degenerate-perturbation}
The entropy production landscape is flat at leading order (Table~\ref{tab:alpha_scan}: $\mathcal{S}/\mathcal{S}_{\max} > 0.92$ across the viable range). The observed configuration may be selected by second-order cross-couplings between parameter sectors, analogous to degenerate perturbation theory. This conjecture requires the multi-parameter Hessian computation listed in the Conclusion; it has not been verified.
\end{remark}

% ===================================================================
\section{Testable Predictions and Falsifiability}
\label{sec:predictions}
% ===================================================================

\subsection{Principal Prediction: Entropy Dominance of the Standard Model}

\begin{conjecture}[Entropy Dominance]
\label{conj:entropy-dominance}
The observed Standard Model parameter set $\vec{p}_{\mathrm{SM}}$ corresponds to a local maximum of the integrated entropy production functional $\mathcal{S}$ within the structurally stable region of parameter space:
\begin{equation}
\mathcal{S}(\mathrm{SM}) > \mathcal{S}(\mathrm{SM}')
\end{equation}
for any varied parameter set $\mathrm{SM}'$ where $|\Delta p_i| \geq 10\%$ for at least one parameter $p_i$, subject to structural stability constraints (nuclear stability, atomic binding, chemical complexity).
\end{conjecture}

\subsection{Concrete Numerical Tests}

\textbf{P1a: Electron-to-Proton Mass Ratio ($m_e/m_p$):} Running stellar evolution codes across a grid of values should reveal that integrated $\Delta S$ is constrained near the observed value by multi-scale compatibility: stellar entropy production is monotonic in $m_e/m_p$ (Section~\ref{sec:stellar_model}), so the observed ratio is selected by atomic and chemical stability bounds rather than a single-scale entropy peak. Significant deviations force stellar engines into inefficient regimes \citep{Adams2019}.

\textbf{P1b: Fine-Structure Constant ($\alpha_{EM}$):} We expect a broad optimization plateau near the observed value (Table~\ref{tab:alpha_scan} confirms $\mathcal{S}/\mathcal{S}_{\max} > 0.92$ for $0.5 \leq \alpha/\alpha_{\text{obs}} \leq 2.0$), with significant decline only beyond $\alpha/\alpha_{\text{obs}} \lesssim 0.5$ or $\gtrsim 3.0$ due to loss of stable chemical bonding, atomic stability, or truncated stellar lifetimes.

\subsection{Semi-Analytic Stellar Model: Quantitative Results}
\label{sec:stellar_model}

As a \emph{pilot study} providing a first quantitative test of the entropy dominance prediction, we construct a semi-analytic stellar model that computes the integrated entropy production $\mathcal{S}_{\text{total}} = \int_0^{t_{\text{MS}}} \dot{S}(t)\,dt \approx E_{\text{nuc}} / T_{\text{eff}}$ as a function of the fine-structure constant $\alpha$ and the electron-to-proton mass ratio $m_e/m_p$. The approximation $\mathcal{S}_{\text{total}} \approx E_{\text{nuc}} / T_{\text{eff}}$ is valid because the effective temperature of a main-sequence star varies by less than a factor of two over its hydrogen-burning lifetime, so that $T_{\text{eff}}$ may be treated as approximately constant for the purpose of this estimate.

\textbf{Model assumptions.} We consider a solar-mass star on the main sequence, modeled as a polytrope with index $n = 3$. The core temperature is determined self-consistently from the balance between nuclear energy generation (dominated by the pp-chain with Gamow tunneling factor) and radiative energy transport (Thomson opacity $\kappa \propto m_e^{-2}$). The Gamow energy for proton-proton fusion is:
\begin{equation}
E_G = 2 m_r c^2 (\pi \alpha)^2 = (\pi \alpha)^2 m_p c^2 \approx 493\,\text{keV}
\end{equation}
where $m_r = m_p/2$ is the reduced mass of the proton-proton system, and the tunneling parameter at solar core temperature is $\tau_\odot = 3(E_G / 4 k_B T_c)^{1/3} \approx 13.5$. The luminosity scales as $L \propto T_c^{\nu}$ where $\nu = \tau/3 - 2/3 \approx 3.8$ for pp-chain conditions. The main-sequence lifetime follows from $t_{\text{MS}} = E_{\text{nuc}} / L$.

\textbf{Note on absolute timescales.} The model computes $t_{\text{MS}}$ from the total nuclear energy reservoir divided by the luminosity, yielding values that differ from detailed stellar evolution codes (which give ${\sim}10\,$Gyr for a solar-mass star) by a factor of ${\sim}7$. This systematic offset arises from the simplified polytropic treatment (no opacity tables, no convective envelope, no metallicity effects). Since all entries in Table~\ref{tab:alpha_scan} are computed with the same model, the \emph{relative} trends---in particular the ratio $\mathcal{S}/\mathcal{S}_{\text{max}}$---are robust, while the absolute timescales should not be compared directly to observed stellar lifetimes.

\textbf{Alpha variation.} Table~\ref{tab:alpha_scan} summarizes the results of varying $\alpha / \alpha_{\text{obs}}$ from 0.25 to 5.0 while holding $m_e/m_p$ and $\alpha_s$ fixed. The strong coupling constant is held constant because its variation would fundamentally alter nuclear physics (binding energies, fusion pathways) in ways that cannot be captured by the polytropic stellar model. The key result is that the integrated entropy production $\mathcal{S}_{\text{total}}$ exhibits a broad maximum near $\alpha / \alpha_{\text{obs}} \approx 1.6$, with the observed value achieving $\mathcal{S}(\alpha_{\text{obs}}) / \mathcal{S}_{\text{max}} = 0.988$.

\begin{table}[ht]
\centering
\caption{Integrated stellar entropy production vs.\ fine-structure constant. The observed value $\alpha_{\text{obs}} = 1/137$ achieves 98.8\% of the unconstrained maximum.}
\label{tab:alpha_scan}
\begin{tabular}{cccccc}
\toprule
$\alpha / \alpha_{\text{obs}}$ & $L / L_\odot$ & $T_{\text{eff}}$ [K] & $t_{\text{MS}}$ [Gyr] & $\mathcal{S}_{\text{total}}$ [J/K] & $\mathcal{S} / \mathcal{S}_{\text{max}}$ \\
\midrule
0.25 & 15.52 & 7026 & 4.67 & $1.247 \times 10^{41}$ & 0.811 \\
0.50 & 4.00 & 6171 & 18.1 & $1.420 \times 10^{41}$ & 0.924 \\
0.75 & 1.77 & 5889 & 41.0 & $1.487 \times 10^{41}$ & 0.968 \\
\textbf{1.00} & \textbf{1.00} & \textbf{5772} & \textbf{72.5} & $\mathbf{1.518 \times 10^{41}}$ & \textbf{0.988} \\
1.25 & 0.63 & 5720 & 114 & $1.532 \times 10^{41}$ & 0.997 \\
1.50 & 0.44 & 5701 & 165 & $1.536 \times 10^{41}$ & 1.000 \\
2.00 & 0.25 & 5712 & 290 & $1.534 \times 10^{41}$ & 0.998 \\
3.00 & 0.11 & 5804 & 640 & $1.508 \times 10^{41}$ & 0.981 \\
5.00 & 0.04 & 6043 & 1813 & $1.450 \times 10^{41}$ & 0.943 \\
\bottomrule
\end{tabular}
\end{table}

\textbf{The plateau problem.} A critical feature of Table~\ref{tab:alpha_scan} is the flatness of the entropy landscape: $\mathcal{S}/\mathcal{S}_{\max} > 0.92$ across the entire range $0.5 \leq \alpha/\alpha_{\text{obs}} \leq 5.0$, and $> 0.96$ for $\alpha/\alpha_{\text{obs}} \in [0.75, 3.0]$. This means that \emph{almost any} value of $\alpha$ within the physically viable range produces near-maximal stellar entropy. The 98.8\% figure therefore does not demonstrate sharp selection of $\alpha_{\text{obs}}$ but rather consistency with a broad plateau. Furthermore, the unconstrained maximum lies at $\alpha/\alpha_{\text{obs}} \approx 1.6$, not at 1.0---the observed value is approximately 37\% below the peak. The framework interprets this as evidence for stability selection from a degenerate manifold rather than sharp optimization: many configurations are near-optimal at first order, and the observed value is selected by second-order cross-couplings between parameter sectors. However, this stability-selection mechanism has not been computed; the Hessian is available only in one dimension. Until the multi-parameter Hessian is evaluated, the 98.8\% figure should be understood as a \emph{necessary but not sufficient} condition for the framework's consistency.

\textbf{Structural stability constraints.} The unconstrained maximum at $\alpha / \alpha_{\text{obs}} \approx 1.6$ must be evaluated against physical viability: for $\alpha \gtrsim 1/85$, stable atoms heavier than iron cease to exist (the fine-structure expansion breaks down), and for $\alpha \gtrsim 1/60$, the triple-alpha resonance shifts out of its viable window \citep{Adams2019}. When these constraints are imposed, the viable parameter window narrows to approximately $0.5 \lesssim \alpha/\alpha_{\text{obs}} \lesssim 2.0$, within which the observed value is near the constrained maximum.

\textbf{Electron mass variation: multi-scale constraints.} \emph{Physical validity bounds must be stated before the model result.} The model's $m_e$ dependence operates through the Thomson opacity $\kappa \propto m_e^{-2}$. However, for sufficiently small $m_e$, the model's assumptions break down: ionization energies scale as $E_I \sim m_e \alpha^2 c^2$, and below $m_e/m_{e,\text{obs}} \lesssim 0.01$, atoms at stellar surface temperatures are fully ionized, rendering the concepts of ``main-sequence star'' and ``stellar opacity'' ill-defined. The Bohr radius $a_0 = \hbar/(m_e \alpha c)$ sets the atomic length scale; if $a_0$ exceeds intermolecular spacings, condensed matter ceases to exist. Chemical bond energies $E_{\text{chem}} \sim m_e \alpha^2 c^2$ must satisfy $E_{\text{chem}} \gg k_B T$ for stable chemistry, yielding a lower bound $m_e \gtrsim 10^{-2}\,m_{e,\text{obs}}$. These constraints impose a hard lower cutoff at approximately $m_e/m_p \gtrsim 10^{-5}$--$10^{-4}$.

\emph{Within} these physical validity bounds, the stellar model reveals a monotonic increase of $\mathcal{S}_{\text{total}}$ with decreasing $m_e$, reflecting the absence of an intrinsic lower bound on $m_e$ at the stellar scale alone. This is a structurally important result: \emph{stellar entropy production alone cannot fix $m_e/m_p$}. The observed value is constrained by the microscopic consistency conditions enumerated above, not by a stellar entropy peak.

\textbf{Two-dimensional scan.} A joint scan over $\alpha$ and $m_e/m_p$ confirms this multi-scale picture quantitatively. The unconstrained 2D maximum lies at $(\alpha/\alpha_{\text{obs}}, m_e/m_{e,\text{obs}}) \approx (1.24, 0.30)$ with $\mathcal{S}_{\text{total}} = 3.22 \times 10^{41}\,\text{J/K}$, yielding $\mathcal{S}(\text{obs})/\mathcal{S}_{\text{max}}^{2D} = 0.47$. The 2D optimum is dominated by the monotonic $m_e$ dependence (lower electron mass $\to$ lower opacity $\to$ longer main-sequence lifetime $\to$ more cumulative entropy). This confirms that \emph{stellar-scale entropy production alone selects only $\alpha$, not $m_e/m_p$}: the observed electron mass requires additional constraints from atomic, chemical, and nuclear physics. The hierarchical selection principle---where MEPP operates within windows defined by lower-scale consistency---is thus not merely a theoretical postulate but is supported by explicit computation.

Within this chemically allowed window, stellar physics operates near maximal entropy production. Thus $m_e/m_p$ is selected by \emph{multi-scale compatibility} rather than by a single-scale MEPP extremum: the tightest constraint comes from the atomic/chemical scale, and larger scales (stellar, galactic) optimize within that permitted range. This hierarchical selection principle is a central structural insight of the FST framework (cf.\ FST-II \citep{FSTII}).

\textbf{Falsifiability assessment.} Among 13 parameter variations tested ($\pm 10\%$, $\pm 20\%$, $\pm 50\%$, $+100\%$ in both $\alpha$ and $m_e/m_p$), 7 yield $\mathcal{S} > \mathcal{S}_{\text{obs}}$ in the unconstrained single-star model. However, the violations in $\alpha$ are small (maximum 1.2\% above observed), and all $m_e$ violations occur at unphysically small electron masses. When structural stability constraints from atomic physics, nuclear stability, and chemical viability are imposed, the number of genuine violations decreases. The 2D scan result ($\mathcal{S}(\text{obs})/\mathcal{S}_{\text{max}}^{2D} = 0.47$) demonstrates that the unconstrained MEPP argument is insufficient: the full FST-I claim requires the hierarchical constraint structure developed above. A comprehensive multi-parameter scan incorporating all relevant physical constraints is in preparation.

\subsection{The Higgs VEV and Cosmological Coupling}

\textbf{P2:} The Higgs VEV is quantitatively related to the cosmological saturation parameter $\Phi_0$. Both values act as order parameters and infrared fixed points \citep{HiggsVEV2024}. A discovered correlation between the scale of electroweak symmetry breaking (246 GeV) and cosmological parameters would support the scale-invariant framework.

\subsection{Neutrino Masses as Optimization Result}

\textbf{P3:} Neutrino masses are non-zero but minimal because they act as efficient ``entropy valves'' in high-density collapse processes.

In core-collapse supernovae, roughly $3 \times 10^{53}$ ergs of gravitational binding energy must be radiated almost entirely by neutrino flux \citep{SupernovaNeutrinos,SupernovaSimulations}. If neutrinos were strictly massless, they would be purely left-handed (right-handed antineutrinos), precluding flavor oscillations and thereby reducing the efficiency of energy transport across different neutrino species during collapse; if too massive, phase-space suppression would throttle the neutrino luminosity and reduce the total energy dissipation rate. The sub-eV mass scale permits flavor oscillations that redistribute energy among all three species, optimizing the transparency of the neutrinosphere and preventing runaway collective oscillations during critical phases like the neutronization burst \citep{NeutrinoOscillations}.

Their specific sub-eV masses are \emph{consistent with} an optimized configuration for high-density entropy extraction. \textbf{Concrete constraint:} The cosmological upper bound on the sum of neutrino masses is $\sum m_\nu < 0.12\,\text{eV}$ (95\% C.L., Planck 2018 \citep{Planck2018}). We observe that the FST-I framework is consistent with masses near the lower boundary of the entropy-efficient window; however, this is a qualitative post-hoc observation, not a quantitative prediction. A precise prediction of the optimal neutrino mass from entropy considerations has not been computed. A precise prediction of the mass hierarchy (normal vs.\ inverted) or the lightest neutrino mass requires a quantitative calculation of the cooling optimization that is not yet available. This is testable with JUNO, DUNE, and future cosmological surveys. Furthermore, the framework predicts an absence of heavily coupled sterile neutrinos that would disrupt the optimized cooling valve.

\subsection{Honest Demarcation}

The framework does not currently predict:
\begin{itemize}
    \item The exact number of fermion generations
    \item The detailed flavor mixing structure (CKM and PMNS matrices)
    \item The specific microscopic nature of dark matter candidates
\end{itemize}

This reflects that MEPP constrains the viable parameter space but does not necessarily dictate every degree of freedom uniquely if multiple configurations yield degenerate payoffs.

\subsection{Strict Falsifiability}

\begin{quote}
\textit{If a theoretical variation of the Standard Model parameters is discovered that significantly increases the integrated entropy production functional $\mathcal{S}$ while simultaneously preserving all structural stability constraints (nuclear stability, atomic binding, stellar longevity, and chemical complexity), then the hypothesis that the Standard Model represents a stability-selected thermodynamic configuration is falsified. The semi-analytic stellar model (Section~\ref{sec:stellar_model}) demonstrates that unconstrained entropy production peaks at $\alpha/\alpha_{\mathrm{obs}} \approx 1.6$, achieving only marginally higher values. The observed fine-structure constant lies within the broad plateau ($\mathcal{S}/\mathcal{S}_{\max} > 0.96$ for $\alpha/\alpha_{\mathrm{obs}} \in [0.75, 3.0]$), consistent with stability selection from a degenerate manifold rather than sharp optimization.}
\end{quote}

\noindent Concretely, a falsifying variation would be a parameter perturbation $\delta p$ in the Standard Model parameter space that (i)~preserves all known stability constraints (vacuum stability, BBN compatibility, stellar ignition), (ii)~increases integrated entropy production $\mathcal{S}[p + \delta p] > \mathcal{S}[p]$, and (iii)~is not excluded by existing experimental bounds. The identification of such a perturbation---or a proof that none exists---constitutes the central open challenge of this program.

\textbf{Concrete quantitative criterion.} To prevent the structural stability constraints from serving as an unfalsifiable escape clause, we commit to the following pre-registered test: when the multi-parameter scan (Conclusion, item~1) is completed over the five core parameters $\{\alpha_{EM}, \alpha_s, G_F, m_e/m_p, m_u/m_d\}$, if more than two of these five parameters lie outside the 80\% plateau of $\mathcal{S}/\mathcal{S}_{\max}$ (i.e., $\mathcal{S}(\text{obs})/\mathcal{S}_{\max} < 0.80$ along the respective single-parameter slice while all other parameters are held at observed values and all structural stability constraints are satisfied), the entropy-dominance conjecture is falsified in its current form. The list of structural stability constraints is fixed at those enumerable from established nuclear and atomic physics: vacuum stability, BBN yield within $2\sigma$ of observed abundances, stellar ignition requiring $T_c > 4 \times 10^6\,$K, and chemical binding energies $E_{\text{chem}} \gg k_B T_{\text{surface}}$. No post-hoc addition of new constraints is permitted.

\subsection{Unified Stability Protocol}\label{sec:stability-protocol}

The falsifiability criteria above require a concrete, reproducible test procedure.
The following protocol operationalizes the resolvent-stability concept for
any system described within the FST framework, with domain-specific
instantiations in FST-I (particle physics), FST-II (chemical networks),
and FST-III (protein folding).

\begin{definition}[FST Stability Test]\label{def:stability-test}
Given a system parameterized by $p \in \mathcal{P}$, the \emph{FST stability
test} consists of three steps:
\begin{enumerate}
  \item \textbf{Perturbation.} Apply a perturbation $\delta p$ to the
        equilibrium configuration $p^\star$.  The nature of $\delta p$ is
        domain-specific:
        \begin{itemize}
          \item FST-I: $\delta p$ represents variations of Standard Model
                parameters (coupling constants, masses, mixing angles).
          \item FST-II: $\delta p$ represents changes in network topology
                or reaction rate constants.
          \item FST-III: $\delta p$ represents conformational changes or
                point mutations in the amino-acid sequence.
        \end{itemize}
  \item \textbf{Primary diagnostic.} Compute the spectral radius
        $\rho(J)$ of the best-response Jacobian
        $J = I - \eta\, H$ at $p^\star$, where $\eta > 0$ is a learning-rate parameter controlling the step size of best-response dynamics (for sufficiently small $\eta$ the spectral-radius criterion $\rho(J)<1$ is equivalent to negative-definiteness of~$H$; see FST-III for an empirical demonstration that the verdict is $\eta$-insensitive over a wide range)\footnote{Sign convention: here $H = -\nabla^2 S$ denotes the \emph{negative} Hessian of the potential, so that $H$ is positive-definite at a local maximum of~$S$. With this convention, $J = I - \eta H$ has eigenvalues $1 - \eta \lambda_i$ with $\lambda_i > 0$, yielding $\rho(J) < 1$ for $0 < \eta < 2/\lambda_{\max}$. The stability verdict (positive-definiteness of~$H$) is independent of~$\eta$; only the convergence rate of the discrete iteration depends on the step size.} and $H$ is the
        stability operator derived from the potential
        (Assumption~\ref{ass:potential-game}).  Exclude gauge directions
        via projection onto the physical subspace.
  \item \textbf{Secondary diagnostic.} Evaluate the entropy production
        rate $\langle\dot{S}\rangle$ as a \emph{correlate}, not a
        definition, of stability.
\end{enumerate}
The system is \emph{resolvent-stable} if and only if $\rho(J) < 1$ for
all non-gauge modes.
\end{definition}

\begin{remark}[Falsifiability]\label{rem:falsifiability-protocol}
The FST stability test yields two testable predictions:
\begin{enumerate}[label=(\roman*)]
  \item Observed equilibrium configurations are resolvent-stable,
        i.e.\ $\rho(J) < 1$ in all non-gauge directions.
  \item Stability and entropy production are positively correlated across
        the accessible parameter space.
\end{enumerate}
A systematic absence of this correlation, or a persistent $\rho(J) > 1$
at an observed equilibrium configuration, would falsify the framework.
For an explicit demonstration in the chemical context, see the
three-species autocatalysis model in FST-II~\cite{FSTII}.
\end{remark}

\subsection{Computational Verification: Fine-Structure Constant}
\label{sec:comp-verification-alpha}

As an independent check of the semi-analytic stellar model (Section~\ref{sec:stellar_model}), we implemented a Gamow-enhanced numerical model that evaluates the integrated entropy production $\mathcal{S}_{\text{total}}$ as a continuous function of $\alpha$ and locates the maximum via gradient-based optimization. Two model variants were compared: a \emph{simple scaling} model that treats the luminosity as a power law $L \propto \alpha^4$ without nuclear physics corrections, and a \emph{Gamow-enhanced} model that incorporates the $\alpha$-dependent Gamow tunneling factor for pp-chain fusion. The results are summarized in Table~\ref{tab:comp-verification-alpha}.

\begin{table}[ht]
\centering
\caption{Computational verification of the entropy-production maximum with respect to the fine-structure constant. The dimensionless coordinate $x_\alpha = \alpha / \alpha_{\text{obs}}$ parameterizes the scan; $x_{\max}$ denotes the location of the entropy maximum; $\mathcal{S}_{\text{obs}} / \mathcal{S}_{\max}$ is the ratio of the observed entropy production to its maximum value; $H(\alpha_{\text{obs}})$ is the Hessian (second derivative) at the observed value. Script: \texttt{stellar\_testbed.py}.}
\label{tab:comp-verification-alpha}
\begin{tabular}{lccc}
\toprule
\textbf{Model} & $x_{\max}$ & $\mathcal{S}_{\text{obs}} / \mathcal{S}_{\max}$ & $H(\alpha_{\text{obs}})$ \\
\midrule
Simple Scaling & 0.658 & 43.3\% & $+6.75$ (minimum) \\
Gamow-Enhanced & 1.605 & 98.75\% & $-0.18$ (maximum) \\
\bottomrule
\end{tabular}
\end{table}

The Gamow-enhanced model places the entropy-production maximum at $x_\alpha = 1.605$, in agreement with the semi-analytic Table~\ref{tab:alpha_scan} result ($x_\alpha \approx 1.5$--$1.6$). The observed fine-structure constant achieves 98.75\% of this maximum. Crucially, the Hessian at $\alpha_{\text{obs}}$ is $H = -0.18$, confirming that the observed value lies on the concave (maximum) side of the entropy landscape. The simple scaling model, which omits the Gamow tunneling dependence on $\alpha$, yields a qualitatively different picture: the maximum shifts to $x_\alpha = 0.658$ and the observed value achieves only 43.3\% of $\mathcal{S}_{\max}$, with a positive Hessian indicating a local minimum. This comparison demonstrates that nuclear physics---specifically the exponential sensitivity of the tunneling rate to $\alpha$---is essential for the MEPP argument: the Gamow factor creates the broad maximum near $\alpha_{\text{obs}}$ that makes the entropy-dominance prediction non-trivial.

The Gamow-enhanced model thus provides a first quantitative test: the observed fine-structure constant lies within 1.25\% of the entropy-production maximum ($\mathcal{S}_{\text{obs}}/\mathcal{S}_{\max} = 0.9875$), with negative curvature $H = -0.18$ confirming a local maximum. While this result is model-dependent (it relies on the polytropic stellar model and the simplified Gamow treatment), the qualitative conclusion---that $\alpha_{\text{obs}}$ sits near a broad entropy maximum when nuclear tunneling is included---is structurally robust.

% ===================================================================
\section{Cosmological Implications and the Arrow of Time}
\label{sec:cosmology}
% ===================================================================

The thermodynamic game governing fundamental-particle sectors unfolds against the dynamically expanding FLRW metric. We note---without claiming explanatory power---that late-time cosmic acceleration is \emph{qualitatively consistent} with entropy maximization: the expansion of phase-space volume maintains the capacity for continued dissipation as localized nuclear fuel is exhausted. This observation is not new (see \citep{EntropyMaximum} for a quantitative discussion of cosmological entropy budgets) and does not constitute a prediction of the cosmological constant from FST-I principles. A quantitative derivation has not been achieved; this remains speculative (Table~\ref{tab:status-claims}).

\subsection{Stability vs.\ Maximization}

The flat entropy landscape (Table~\ref{tab:alpha_scan}: $\mathcal{S}/\mathcal{S}_{\max} > 0.92$ across the viable range) suggests that the fine-tuning problem may be better framed as a stability problem than a maximization puzzle. Many configurations yield comparable first-order entropy production; the observed configuration may be selected by second-order cross-couplings between parameter sectors. This ``stability selection'' hypothesis provides an alternative to both the landscape approach and anthropic reasoning, but it remains a \emph{structural conjecture}: the full Hessian matrix $H_{ij} = \partial^2 \mathcal{S}/\partial p_i \partial p_j$ has not been computed. The only available datum is the 1D curvature $H(\alpha_{\text{obs}}) = -0.18$ from the Gamow-enhanced model (Table~\ref{tab:comp-verification-alpha}). Whether the multi-parameter Hessian is negative-definite in all non-gauge directions---the necessary and sufficient condition for stability selection---is an open question whose resolution requires the multi-parameter scan listed as priority item~2 in the Conclusion.

\subsection{Status of Claims}

Table~\ref{tab:status-claims} classifies the central claims of this paper by epistemic status.

\begin{table}[htbp]
\centering
\caption{Status of claims. \textsc{Established}: proven or widely accepted results cited here; \textsc{Conjectured}: new but testable claims; \textsc{Speculative}: framework-level proposals without direct empirical test.}
\label{tab:status-claims}
\small
\begin{tabularx}{\textwidth}{@{}>{\raggedright\arraybackslash}p{4.0cm}>{\raggedright\arraybackslash}p{2.7cm}Y@{}}
\toprule
\textbf{Claim} & \textbf{Status} & \textbf{Evidence / Reference} \\
\midrule
Nash equilibrium / mean field game theory & \textsc{Established} & Nash (1950); Lasry \& Lions (2007) \\
MEPP as heuristic organizing principle & \textsc{Established} & Martyushev \& Seleznev (2006); Dyke \& Kleidon (2010) \\
Born rule via envariance & \textsc{Established} & Zurek (2003, 2005); contested scope \\
Diproton catastrophe \& triple-alpha sensitivity & \textsc{Established} & Adams (2019) \\
FEP--England duality (Conjecture~\ref{lem:fep-england}) & \textsc{Conjectured} & Heuristic from Crooks theorem + FEP; no rigorous proof published \\
Quantum--strategy correspondence (commutative specialization) & \textsc{Conjectured} & Sec.~\ref{sec:foundations}; diagonal reduction of Lindblad-OT to Lasry--Lions conjectured here; requires independent verification; coherences require noncommutative extension \\
SM parameters as Nash equilibrium of entropy game & \textsc{Conjectured} & Conjecture~\ref{conj:entropy-dominance}; pilot study only \\
Stellar entropy peaks near $\alpha_{\text{obs}}$ (98.8\%) & \textsc{Conjectured} & Table~\ref{tab:alpha_scan}; single-channel, single-star model \\
Resolvent stability as fine-tuning explanation & \textsc{Conjectured} & Sec.~\ref{sec:cosmology}; structural analogy, no Hessian computed \\
Symmetry breaking as phase transition between equilibria & \textsc{Conjectured} & Sec.~\ref{sec:symmetry}; reinterpretation of known physics \\
Neutrino masses as entropy-valve optimization & \textsc{Speculative} & P3; no quantitative mass prediction \\
Higgs VEV--cosmological saturation connection & \textsc{Speculative} & P2; no derivation of $\Lambda$ from MEPP \\
Cosmological acceleration as entropy strategy & \textsc{Speculative} & Sec.~\ref{sec:cosmology}; heuristic only \\
\bottomrule
\end{tabularx}
\end{table}

% ===================================================================
\section{Conclusion}
\label{sec:conclusion}
% ===================================================================

This work proposes a unifying \emph{interpretive framework} for fundamental physics based on a single organizing principle: thermodynamic optimization under strategic interaction. Within this framework, the Standard Model is not viewed as an arbitrary collection of parameters, but as potentially consistent with the constrained equilibrium outcome of a multi-scale optimization game. We emphasize that FST-I is, at this stage, a structured hypothesis rather than a predictive theory: it generates no quantitative predictions beyond established physics and derives no parameter values from first principles. Its contribution is a conceptual vocabulary that organizes known parameter sensitivities within a game-theoretic architecture.

At the particle level, quantum mechanics is reinterpreted as a mixed-strategy equilibrium, with the path integral acting as a Nash selector that suppresses non-optimal strategies through destructive interference. The Born rule is recast via envariance as the unique equilibrium distribution---a conceptual reinterpretation, not a new derivation. Symmetry breaking is reinterpreted as a phase transition between equilibria, where changing external conditions destabilize symmetric configurations and drive the system toward entropy-producing attractors. Fine-tuning, in this view, is not a mystery to be explained away, but the expected signature of an optimized solution.

Crucially, this framework does not modify the formal structure of established physics. No new particles, forces, or dynamical laws are introduced. Instead, the contribution lies in identifying a common optimization logic that organizes known phenomena across scales. The Standard Model parameters are interpreted as a Nash-optimal toolkit for entropy production in a universe governed by thermodynamic constraints.

The approach includes explicit falsifiability criteria (Section~\ref{sec:predictions}), including a pre-registered quantitative criterion for the planned multi-parameter scan. A first semi-analytic pilot study demonstrates that the observed fine-structure constant lies within a broad plateau of near-maximal stellar entropy production ($\mathcal{S}/\mathcal{S}_{\max} = 0.988$, but with $\mathcal{S}/\mathcal{S}_{\max} > 0.92$ across the entire viable range)---a consistency check, not a decisive test. The flatness of this plateau means that the 1D result, while necessary for the framework's viability, does not demonstrate sharp parameter selection. A two-dimensional scan reveals $\mathcal{S}/\mathcal{S}_{\max}^{2D} = 0.47$, demonstrating that single-scale stellar MEPP is falsified as a stand-alone selection principle: without hierarchical multi-scale constraints from atomic and chemical physics, the entropy dominance conjecture fails already in two dimensions. The full multi-scale constraint hierarchy is therefore not an optional refinement but a necessary structural component of the framework. Extending this pilot study to the full Standard Model parameter space remains the central computational challenge.

This paper represents the first step of a broader research program. The immediate next steps, in order of priority, are:
\begin{enumerate}
    \item \textbf{Multi-parameter stellar scan (critical):} Extend the semi-analytic model to vary $\alpha_s$, $G_F$, and $m_u/m_d$ simultaneously, incorporating nuclear stability boundaries from \citep{Adams2019}. This is the most direct route to testing the entropy-dominance conjecture beyond the single-parameter results presented here.
    \item \textbf{Hessian analysis (critical):} Compute the second-order stability matrix $\partial^2 \mathcal{S}/\partial p_i \partial p_j$ at the observed parameter point to test the resolvent-stability conjecture quantitatively. A negative-definite Hessian in all non-gauge directions would provide the strongest evidence for the framework.
    \item \textbf{Multi-channel entropy:} Include gravitational (structure formation), chemical (molecular complexity), and nuclear (nucleosynthesis yield) entropy channels beyond stellar radiation.
    \item \textbf{Full stellar evolution codes:} Replace the polytropic model with MESA-type stellar evolution simulations using non-standard nuclear parameters.
\end{enumerate}

If the computational program outlined above succeeds in demonstrating that the observed parameters occupy a multi-dimensional stability maximum, the picture would suggest that physical laws at different scales share a common optimization logic rather than being independently constructed. In that case, fine-tuning would cease to be a problem and become a diagnostic: a trace of equilibrium in a thermodynamic game spanning from the quantum vacuum to the largest cosmic structures. Until then, FST-I remains a promising but unvalidated interpretive framework.

% ===================================================================
% References
% ===================================================================

\bibliographystyle{plainnat}
{\small
\setlength{\parskip}{0pt}
\begin{thebibliography}{99}
\setlength{\itemsep}{0pt}

\bibitem[Aghanim et al.(2020)]{Planck2018}
Aghanim, N. et al. (Planck Collaboration) (2020).
\newblock Planck 2018 results. VI. Cosmological parameters.
\newblock \textit{Astronomy \& Astrophysics}, 641, A6.
\newblock arXiv:1807.06209.

\bibitem[Adams(2019)]{Adams2019}
Adams, F.~C. (2019).
\newblock The degree of fine-tuning in our universe---and others.
\newblock \textit{Physics Reports}, 807, 1--111.
\newblock arXiv:1902.03928.

\bibitem[Crooks(1999)]{Crooks1999}
Crooks, G.~E. (1999).
\newblock Entropy production fluctuation theorem and the nonequilibrium work relation for free energy differences.
\newblock \textit{Physical Review E}, 60, 2721--2726.

\bibitem[England(2013)]{England2013}
England, J.~L. (2013).
\newblock Statistical physics of self-replication.
\newblock \textit{The Journal of Chemical Physics}, 139, 121923.
\newblock DOI: 10.1063/1.4818538.

\bibitem[Dewar(2003)]{Dewar2003}
Dewar, R.~C. (2003).
\newblock Information theory explanation of the fluctuation theorem, maximum entropy production and self-organized criticality in non-equilibrium stationary states.
\newblock \textit{Journal of Physics A: Mathematical and General}, 36(3), 631.

\bibitem[Djete and Touzi(2026)]{DjeteTouzi2026}
M.~F.~Djete and N.~Touzi (2026).
\newblock On Approximate Nash Equilibria in Mean Field Games.
\newblock arXiv:2601.20910.

\bibitem[Grinstein \& Linsker(2007)]{GrinsteinLinsker2007}
Grinstein, G. \& Linsker, R. (2007).
\newblock Comments on a derivation and application of the ``maximum entropy production'' principle.
\newblock \textit{Journal of Physics A: Mathematical and Theoretical}, 40(31), 9717.

\bibitem[Bousso \& Polchinski(2000)]{Bousso2000}
Bousso, R. \& Polchinski, J. (2000).
\newblock Quantization of four-form fluxes and dynamical neutralization of the cosmological constant.
\newblock \textit{JHEP}, 06, 006.

\bibitem[Valentini \& Westman(2005)]{ValentiniWestman2005}
Valentini, A. \& Westman, H. (2005).
\newblock Dynamical origin of quantum probabilities.
\newblock \textit{Proc.\ R.\ Soc.\ A}, 461(2053), 253--272. arXiv:quant-ph/0403034.

\bibitem[Friston(2010)]{Friston2010}
Friston, K. (2010).
\newblock The free-energy principle: a unified brain theory?
\newblock \textit{Nature Reviews Neuroscience}, 11, 127--138.

\bibitem[Friston et al.(2006)]{FristonFreeEnergy}
Friston, K., Kilner, J. \& Harrison, L. (2006).
\newblock A free energy principle for the brain.
\newblock \textit{Journal of Physiology -- Paris}, 100, 70--87.

\bibitem[Kleidon \& Lorenz(2004)]{MEPP2004}
Kleidon, A. \& Lorenz, R.~D. (eds.) (2004).
\newblock Non-equilibrium Thermodynamics and the Production of Entropy: Life, Earth, and Beyond.
\newblock \textit{Springer}.

\bibitem[Dyke \& Kleidon(2010)]{MEPPFoundations}
Dyke, J. \& Kleidon, A. (2010).
\newblock The Maximum Entropy Production Principle: Its Theoretical Foundations and Applications to the Earth System.
\newblock \textit{Entropy}, 12, 613--630.
\newblock DOI: 10.3390/e12030613.

\bibitem[Abel(2023)]{Abel2023}
Abel, C.~R. (2023).
\newblock The quantum foundations of utility and value.
\newblock \textit{Phil.\ Trans.\ R.\ Soc.\ A}, 381, 20220286.

% Maity (2024) removed from active citations; retained in working bibliography for reference.
% Original: arXiv:2402.13395 (not peer-reviewed, non-physics venue).

\bibitem[Lindblad(1976)]{Lindblad1976}
Lindblad, G. (1976).
\newblock On the generators of quantum dynamical semigroups.
\newblock \textit{Communications in Mathematical Physics}, 48, 119--130.

\bibitem[Martyushev \& Seleznev(2006)]{MartyushevSeleznev2006}
Martyushev, L.~M. \& Seleznev, V.~D. (2006).
\newblock Maximum entropy production principle in physics, chemistry and biology.
\newblock \textit{Physics Reports}, 426, 1--45.

\bibitem[Martyushev(2010)]{Martyushev2010}
Martyushev, L.~M. (2010).
\newblock The maximum entropy production principle: two basic questions.
\newblock \textit{Phil. Trans. R. Soc. B}, 365(1545), 1333--1334.

\bibitem[Lasry and Lions(2006)]{LasryLions2006}
J.-M.~Lasry and P.-L.~Lions (2006).
\newblock Jeux \`a champ moyen.\ I -- Le cas stationnaire.
\newblock \textit{Comptes Rendus Math\'ematique}, 343(9), 619--625.
\newblock DOI: 10.1016/j.crma.2006.09.019.

\bibitem[Lasry \& Lions(2007)]{LasryLions2007}
Lasry, J.-M. \& Lions, P.-L. (2007).
\newblock Mean field games.
\newblock \textit{Japanese Journal of Mathematics}, 2(1), 229--260.

\bibitem[Padmanabhan(2010)]{Padmanabhan2010}
Padmanabhan, T. (2010).
\newblock Thermodynamical aspects of gravity: New insights.
\newblock \textit{Rep. Prog. Phys.}, 73, 046901.
\newblock DOI: 10.1088/0034-4885/73/4/046901.

\bibitem[Feynman \& Hibbs(1965)]{PathIntegral}
Feynman, R.~P. \& Hibbs, A.~R. (1965).
\newblock \textit{Quantum Mechanics and Path Integrals}.
\newblock McGraw-Hill.

\bibitem[Kleinert(2009)]{PathIntegralOptimization}
Kleinert, H. (2009).
\newblock \textit{Path Integrals in Quantum Mechanics, Statistics, Polymer Physics, and Financial Markets} (5th ed.).
\newblock World Scientific.

% Pramanik (2024) removed from active citations; retained in working bibliography for reference.
% Original: Ada Academica (non-standard venue, not indexed in Scopus/WoS).

\bibitem[Busemeyer \& Bruza(2012)]{BusemeyerBruza2012}
Busemeyer, J.~R. \& Bruza, P.~D. (2012).
\newblock \textit{Quantum Models of Cognition and Decision}.
\newblock Cambridge University Press.

\bibitem[Eisert et al.(1999)]{QuantumOptimization}
Eisert, J., Wilkens, M. \& Lewenstein, M. (1999).
\newblock Quantum games and quantum strategies.
\newblock \textit{Physical Review Letters}, 83(15), 3077--3080.

\bibitem[Chakrabarti et al.(2024)]{HiggsVEV2024}
Chakrabarti, S., Anagha V., Ganesh, S. \& Menon, V. (2024).
\newblock A cosmological reconstruction of the {H}iggs vacuum expectation value.
\newblock \textit{Eur.\ Phys.\ J.\ C}, 84, 1250. arXiv:2409.15962.

\bibitem[Schlosshauer(2007)]{MeasurementClassical}
Schlosshauer, M. (2007).
\newblock \textit{Decoherence and the Quantum-to-Classical Transition}.
\newblock Springer.

\bibitem[Jaynes(1957)]{Jaynes1957}
Jaynes, E.~T. (1957).
\newblock Information theory and statistical mechanics.
\newblock \textit{Physical Review}, 106(4), 620--630.

\bibitem[Janka(2012)]{SupernovaNeutrinos}
Janka, H.-T. (2012).
\newblock Explosion mechanisms of core-collapse supernovae.
\newblock \textit{Annual Review of Nuclear and Particle Science}, 62, 407--451.

\bibitem[Burrows(2013)]{SupernovaSimulations}
Burrows, A. (2013).
\newblock Colloquium: Perspectives on core-collapse supernova theory.
\newblock \textit{Reviews of Modern Physics}, 85(1), 245--261.

\bibitem[Johns et al.(2025)]{NeutrinoOscillations}
Johns, L., Richers, S. \& Wu, M.-R. (2025).
\newblock Neutrino oscillations in core-collapse supernovae and neutron star mergers.
\newblock \textit{Ann.\ Rev.\ Nucl.\ Part.\ Sci.}, 75, 399. arXiv:2503.05959.

\bibitem[Smolin(1992)]{Smolin1992}
Smolin, L. (1992).
\newblock Did the Universe evolve?
\newblock \textit{Classical and Quantum Gravity}, 9, 173--191.

\bibitem[Susskind(2005)]{Susskind2005}
Susskind, L. (2005).
\newblock \textit{The Cosmic Landscape}.
\newblock Little, Brown.

\bibitem[Tegmark et al.(2006)]{Tegmark2006}
Tegmark, M. et al. (2006).
\newblock Dimensionless constants, cosmology, and other dark matters.
\newblock \textit{Phys. Rev. D}, 73, 023505.

\bibitem['t~Hooft(2016)]{tHooft2016}
't~Hooft, G. (2016).
\newblock \textit{The Cellular Automaton Interpretation of Quantum Mechanics}.
\newblock Springer, Fundamental Theories of Physics, Vol.~185.

\bibitem[Busch et al.(2007)]{HeisenbergCovariant}
Busch, P., Heinonen, T. \& Lahti, P. (2007).
\newblock Heisenberg's uncertainty principle.
\newblock \textit{Physics Reports}, 452(6), 155--176.

\bibitem[Peskin \& Schroeder(1995)]{HiggsVEV}
Peskin, M.~E. \& Schroeder, D.~V. (1995).
\newblock \textit{An Introduction to Quantum Field Theory}.
\newblock Addison-Wesley. (Ch. 20, electroweak symmetry breaking).

\bibitem[Polettini(2013)]{Polettini2013}
Polettini, M. (2013).
\newblock Fact-checking Ziegler's maximum entropy production principle beyond the linear regime and towards steady states.
\newblock \textit{Entropy}, 15(7), 2570--2584.
\newblock DOI: 10.3390/e15072570.

\bibitem[Verlinde(2011)]{Verlinde2011}
Verlinde, E. (2011).
\newblock On the origin of gravity and the laws of Newton.
\newblock \textit{JHEP}, 04, 029.
\newblock DOI: 10.1007/JHEP04(2011)029.

\bibitem[Weinberg(1987)]{Weinberg1987}
Weinberg, S. (1987).
\newblock Anthropic bound on the cosmological constant.
\newblock \textit{Phys. Rev. Lett.}, 59, 2607.

\bibitem[Egan \& Lineweaver(2010)]{EntropyMaximum}
Egan, C.~A. \& Lineweaver, C.~H. (2010).
\newblock A larger estimate of the entropy of the universe.
\newblock \textit{ApJ}, 710, 1825.

\bibitem[Zurek(2003)]{Zurek2003}
Zurek, W.~H. (2003).
\newblock Decoherence, einselection, and the quantum origins of the classical.
\newblock \textit{Rev. Mod. Phys.}, 75, 715--775.

\bibitem[Zurek(2005)]{Zurek2005}
Zurek, W.~H. (2005).
\newblock Probabilities from entanglement, Born's rule from envariance.
\newblock \textit{Physical Review A}, 71, 052105. arXiv:quant-ph/0405161.
\newblock DOI: 10.1103/PhysRevA.71.052105.

\bibitem[Schlosshauer \& Fine(2005)]{SchlossFine2005}
Schlosshauer, M. \& Fine, A. (2005).
\newblock On Zurek's derivation of the Born rule.
\newblock \textit{Foundations of Physics}, 35, 197--213. arXiv:quant-ph/0312058.

\bibitem[Geiger(2026b)]{FSTII}
Geiger, L. (2026).
\newblock Functional Stability Theory II: Chemical Stability and Autocatalytic Selection.
\newblock \textit{Zenodo preprint}, 2026.
\newblock DOI: 10.5281/zenodo.20130563.

\bibitem[Geiger(2026c)]{GeigerRH2026c}
Geiger, L. (2026).
\newblock The Riemann Hypothesis Part III: The Analytical Rank-Finite Certificate.
\newblock \textit{Zenodo preprint}, March 2026.
\newblock Series Concept DOI: 10.5281/zenodo.19035640.

\bibitem[Geiger(2026e)]{GeigerRH2026b}
Geiger, L. (2026).
\newblock Even Dominance of the Weil Quadratic Form: Spectral Analysis, Computer-Assisted Proof, and the Resolvent Gap Theorem.
\newblock \textit{Zenodo preprint}, March 2026.
\newblock Concept DOI: 10.5281/zenodo.19035640.

\bibitem[Geiger(2026d)]{FST-Unified2026}
Geiger, L. (2026).
\newblock FST Spectrum Duality: The Renormalized Free Energy Principle as Physical Instantiation of Functional Stability Theory.
\newblock \textit{Zenodo preprint}, 2026.
\newblock DOI: 10.5281/zenodo.19036190.

\bibitem[Geiger(2026f)]{GeigerCRM2026}
Geiger, L. (2026).
\newblock The Curvature Relaxation Model: A Four-Paper Program for Geometric Cosmology Without the Dark Sector.
\newblock \textit{Zenodo preprint}, 2026.
\newblock Concept DOI: 10.5281/zenodo.18728935.

\bibitem[Nash(1950)]{Nash1950}
Nash, J.~F. (1950).
\newblock Equilibrium points in $n$-person games.
\newblock \textit{Proceedings of the National Academy of Sciences}, 36(1), 48--49.

\bibitem[Paltridge(1975)]{Paltridge1975}
Paltridge, G.~W. (1975).
\newblock Global dynamics and climate---a system of minimum entropy exchange.
\newblock \textit{Quarterly Journal of the Royal Meteorological Society}, 101(429), 475--484.

\bibitem[Prigogine(1967)]{Prigogine1967}
Prigogine, I. (1967).
\newblock \textit{Introduction to Thermodynamics of Irreversible Processes} (3rd ed.).
\newblock Wiley Interscience.

\bibitem[Sawada et al.(2025)]{Sawada2025}
Sawada, Y., Daigaku, Y. \& Toma, K. (2025).
\newblock Maximum entropy production principle of thermodynamics for the birth and evolution of life.
\newblock \textit{Entropy}, 27(4), 449.
\newblock DOI: 10.3390/e27040449.

\end{thebibliography}

\medskip
\noindent\textbf{AI Disclosure.} This work was assisted by Claude (Anthropic) for literature synthesis, mathematical verification, and text drafting. All scientific claims, original arguments, and theoretical contributions are the sole responsibility of the human author. No AI system is listed as co-author.
}

\end{document}
